\section{Preliminaries}\label{Section2}

Throughout this section,~$G$ denotes any compact Lie groups, and $X, Y, \dots$ any f\/inite $CW$-complexes unless
otherwise specif\/ied.

\vspace{-1mm}

\subsection[(Equivariant) $K$-theory of compact connected Lie groups]{(Equivariant)
$\boldsymbol{K}$-theory of compact connected Lie groups}
%\label{K_GG}

Recall that the functor $K^{-1}$ is represented by $U(\infty):=\lim\limits_{n\to\infty}U(n)$, i.e.\
for any $X$, $K^{-1}(X)$ is the abelian group of homotopy classes of maps $[X, U(\infty)]$.
Such a~description of $K^{-1}$ leads to the following
\begin{Definition}
\label{betamap}
Let $\delta: R(G)\to K^{-1}(G)$ be the group homomorphism which sends any complex~$G$-representation $\rho$ to the
homotopy class of $i\circ \rho$, where $\rho$ is regarded as a~continuous map from~$G$ to $U(n)$ and $i:
U(n)\hookrightarrow U(\infty)$ is the standard inclusion.
\end{Definition}

{\samepage In fact, any element in $K^{-q}(X)$ can be represented by a~complex of vector bundles on $X\times \mathbb{R}^q$, which
is exact outside $X\times\{0\}$ (cf.~\cite{At}).
This gives another interpretation of $\delta(\rho)$, as shown in the following proposition,
which we f\/ind useful for our
exposition.}

\begin{Proposition}
\label{-1_class}
If~$V$ is the underlying complex vector space of $\rho$, then $\delta(\rho)$ is represented by
\begin{gather*}
0\to G\times\mathbb{R}\times V \to G\times\mathbb{R}\times V\to 0,
\\
(g, t, v) \ \mapsto\  (g, t, -t\rho(g)v) \qquad\text{if} \ \ t\geq 0,
\\
(g, t, v) \ \mapsto\  (g, t, tv) \qquad\text{if}\ \  t\leq 0.
\end{gather*}
\end{Proposition}

\begin{Proposition}[\cite{Ho}]\label{noneqderivation}\quad
\begin{enumerate}\itemsep=0pt
\item[$1.$]
$\delta$ is a~derivation of $R(G)$ taking values in $K^{-1}(G)$ regarded as an $R(G)$-module whose module structure is
given by the augmentation map.
In other words, $\delta$ is a~group homomorphism from $R(G)$ to $K^{-1}(G)$ satisfying
\begin{gather}
\label{leibniz}
\delta(\rho_1\otimes\rho_2)=\dim(\rho_1)\delta(\rho_2)+\dim(\rho_2)\delta(\rho_1).
\end{gather}
\item[$2.$]
If $I(G)$ is the augmentation ideal of $R(G)$, then $\delta(I(G)^2)=0$.
\end{enumerate}
\end{Proposition}

The main results of~\cite{Ho} are stated in the following
\begin{Theorem}
\label{hodgkin}
Let~$G$ be a~compact connected Lie group with torsion-free fundamental group.
Then
\begin{enumerate}\itemsep=0pt
\item[$1.$]
$K^*(G)$ is torsion-free.
\item[$2.$]
Let $J(G):=I(G)/I(G)^2$.
Then the map $\widetilde{\delta}: J(G)\to K^{-1}(G)$ induced by $\delta$ is well-defined, and
$K^*(G)=\bigwedge(\text{\rm Im}(\widetilde{\delta}))$.
\item[$3.$]
In particular, if~$G$ is compact, connected and simply-connected of rank $l$, then
$K^*(G)=\bigwedge_{\mathbb{Z}}(\delta(\rho_1), \dots, \delta(\rho_l))$, where $\rho_1, \dots, \rho_l$ are the
fundamental representations.
\end{enumerate}
\end{Theorem}

Viewing~$G$ as a~$G$-space via the adjoint action, one may consider the equivariant $K$-theo\-ry~$K_G^*(G)$.
Let $\Omega^*_{R(G)/\mathbb{Z}}$ be the ring of Grothendieck dif\/ferentials of $R(G)$ over $\mathbb{Z}$, i.e.\
the exterior algebra over $R(G)$ of the module of K\"ahler dif\/ferentials of $R(G)$ over $\mathbb{Z}$ (cf.~\cite{BZ}).
\begin{Definition}
\label{defdeltag}
Let $\varphi: \Omega_{R(G)/\mathbb{Z}}^*\to K_G^*(G)$ be the $R(G)$-algebra homomorphism def\/ined by the following
\begin{enumerate}\itemsep=0pt
\item[1)]
$\varphi(\rho_V):=[G\times V]\in K_G^*(G)$, where~$G$ acts on $G\times V$ by $g_0\cdot(g_1, v)=(g_0g_1g_0^{-1},
\rho_V(g_0)v)$,
\item[2)]
$\varphi(d\rho_V)\in K_G^{-1}(G)$ is the complex of vector bundles in Def\/inition~\ref{-1_class} where~$G$ acts on
$G\times \mathbb{R}\times V$ by $g_0\cdot(g_1, t, v)=(g_0g_1g_0^{-1}, t, \rho_V(g_0)v)$,
\end{enumerate}
We also def\/ine $\delta_G: R(G)\to K^{-1}_G (G)$ by $\delta_G(\rho_V):=\varphi(d\rho_V)$.
\end{Definition}
\begin{Remark}
The def\/inition of $\delta_G(\rho_V)$ given in~\cite{BZ}, where the middle map of the complex of vector bundles is
def\/ined to be $(g, t, v)\mapsto(g, t, t\rho_V(g)v)$ for all $t\in\mathbb{R}$, is incorrect, as $\delta_G(\rho_V)$ so
def\/ined is actually 0 in $K_G^{-1}(G)$.
The def\/inition given in Def\/inition~\ref{defdeltag} is the correction made by Brylinski and relayed to us by one of the
referees.
\end{Remark}
\begin{Theorem}[\cite{BZ}]\label{eqderivation}\quad
\begin{enumerate}\itemsep=0pt
\item[$1.$]
$\delta_G$ is a~derivation of $R(G)$ taking values in the $R(G)$-module $K^{-1}_G(G)$, i.e.\
$\delta_G$ satisfies
\begin{gather}
\label{eqleibniz}
\delta_G(\rho_1\otimes\rho_2)=\rho_1\cdot\delta_G(\rho_2)+\rho_2\cdot\delta_G(\rho_1).
\end{gather}
\item[$2.$]
Let~$G$ be a~compact connected Lie group with torsion-free fundamental group.
Then $\varphi$ is an $R(G)$-algebra isomorphism.
\end{enumerate}
\end{Theorem}

\begin{Remark}\label{weakequivformality}\quad
\begin{enumerate}\itemsep=0pt
\item
Although the def\/inition of $\delta_G(\rho_V)$ given by Brylinski--Zhang in~\cite{BZ} is incorrect, their proof of
Theorem~\ref{eqderivation} can be easily corrected by using the correct def\/inition as in Def\/inition~\ref{defdeltag},
which does not af\/fect the validity of the rest of their arguments, and Theorem~\ref{eqderivation} still stands.
\item
In~\cite{HL}, a~$G$-space $X$ is def\/ined to be \emph{weakly equivariantly formal} if the map
$K^*_G(X)\otimes_{R(G)}\mathbb{Z}\to K^*(X)$ induced by the forgetful map is a~ring isomorphism, where $\mathbb{Z}$ is
viewed as an $R(G)$-module through the augmentation homomorphism.
Theorem~\ref{hodgkin} and~\ref{eqderivation} imply that~$G$ is weakly equivariantly formal if it is connected with
torsion-free fundamental group.
We will make use of this property in computing the equivariant $KR$-theory of~$G$.
\item
Let $f: K_G^*(G)\to K^*(G)$ be the forgetful map.
Note that $f(\varphi(\rho))=\dim(\rho)$ and $f(\varphi(d\rho))=\delta(\rho)$.
Applying $f$ to equation~\eqref{eqleibniz} in Theorem~\ref{eqderivation}, we get equation~\eqref{leibniz} in
Proposition~\ref{noneqderivation}.
\item
$I(G)$, being a~prime ideal in $R(G)$, can be thought of as an element in $\Spec R(G)$, and
$K^*(G)\cong\bigwedge_{\mathbb{Z}}T^*_{I(G)}\Spec R(G)$, $K_G^*(G)\cong\bigwedge_{R(G)}T^*_{I(G)}\Spec R(G)$.
\end{enumerate}
\end{Remark}
\subsection{Real representation rings}
This section is an elaboration of the part on Real representation rings in~\cite{AS} and~\cite{Se}.
Since the results in this subsection can be readily generalized from those results concerning the special case
where $\sigma_G$ is trivial, we omit most of the proofs and refer the reader to any standard text on representation theory of
Lie groups, e.g.~\cite{B}
and~\cite{BtD}.
\begin{Definition}
A Real Lie group is a~pair $(G,\sigma_G)$ where~$G$ is a~Lie group and $\sigma_G$ a~Lie group involution on it.
A \emph{Real representation}
$V$ of a~Real Lie group $(G,\sigma_G)$ is a~f\/inite-dimensional complex representation of~$G$ equipped with an anti-linear involution~$\sigma_V$ such
that $\sigma_V(gv)=\sigma_G(g)\sigma_V(v)$.
Let $\mathcal{R}\text{ep}_\mathbb{R}(G,
\sigma_G)$ be the category of Real representations of $(G, \sigma_G)$.
A morphism between~$V$ and $W\in \mathcal{R}
ep_\mathbb{R}(G,
\sigma_G)$ is a~linear transformation from~$V$ to $W$ which commute with~$G$ and respect~$\sigma_V$ and $\sigma_W$.
We denote $\Mor
(V, W)$ by $\Hom_G(V, W)^{(\sigma_V, \sigma_W)}$.
An irreducible Real representation is an irreducible object in $\mathcal{R}\text{ep}_\mathbb{R}(G,
\sigma_G)$.
The Real representation ring of $(G, \sigma_G)$, denoted by $RR(G, \sigma_G)$, is the Grothendieck group of $\mathcal{R}
ep_\mathbb{R}(G,
\sigma_G)$, with multiplication being tensor product over $\mathbb{C}$.
Sometimes we will omit the notation $\sigma_G$ if there is no danger of confusion about the Lie group involution.
\end{Definition}
\begin{Remark}
Let~$V$ be an irreducible Real representation of~$G$.
Then $\Hom_G(V, V)^{\sigma_V}$ must be a~real associative division algebra which, according to Frobenius' theorem, is isomorphic to either
$\mathbb{R}$, $\mathbb{C}$ or $\mathbb{H}$.
Following the convention of~\cite{AS}, we call $\Hom_G(V, V)^{\sigma_V}$ the \emph{commuting field} of~$V$.
\end{Remark}

\begin{Definition}
If~$V$ is an irreducible Real representation of~$G$, then we say~$V$ is of \emph{real, complex or quaternionic type}
according as the commuting f\/ield is isomorphic to $\mathbb{R}$, $\mathbb{C}$ or $\mathbb{H}$.
Let $RR(G, \mathbb{F})$ be the abelian group generated by the isomorphism classes of irreducible Real representations
with $\mathbb{F}$ as the commuting f\/ield.
\end{Definition}
\begin{Remark}
$RR(G)\cong RR(G, \mathbb{R})\oplus RR(G, \mathbb{C})\oplus RR(G, \mathbb{H})$ as abelian groups.
\end{Remark}
\begin{Definition}
Let~$V$ be a~$G$-representation.
We use $\sigma_G^*V$ to denote the~$G$-representation with the same underlying vector space where the~$G$-action is twisted by
$\sigma_G$, i.e.\
$\rho_{\sigma_G^*V}
(g)v=\rho_V(\sigma_G(g))v$.
We will use $\overline{\sigma_G^*}$ to denote the map on $R(G)$ def\/ined by $[V]\mapsto [\sigma_G^*\overline{V}]$.
\end{Definition}

\begin{Proposition}
%\label{converserealselfconjugate}
If~$V$ is a~complex~$G$-representation, and there exists $f\in\Hom_G(V,\sigma_G^*\overline{V})$
such that $f^2=\Id_V$, then~$V$ is a~Real representation of~$G$ with $f$ as the anti-linear involution~$\sigma_V$.
\end{Proposition}
\begin{proof}
If $f\in\Hom_G(V,
\sigma_G^*\overline{V})$, then it is anti-linear on~$V$ and $f(gv)=
\sigma_G(g)f(v)$ for $g\in G$ and $v\in V$.
The assumption that $f^2=\Id
_V$ just says that $f$ is an involution.
So~$V$ together with $\sigma_V=f$ is a~Real representation of~$G$.
\end{proof}
\begin{Proposition}
%\label{fsindicator}
Let~$V$ be an irreducible complex representation of~$G$ and suppose that $V\cong
\sigma_G^*\overline{V}$.
Let $f\in\Hom_G(V,
\sigma_G^*\overline{V})$.
Then
\begin{enumerate}\itemsep=0pt
\item[$1.$]
$f^2=k\Id_V$ for some $k\in\mathbb{R}$.
\item[$2.$]
There exists $g\in\Hom_G(V,
\sigma_G^*\overline{V})$ such that $g^2=\Id_V$ or $g^2=-\Id_V$.
\end{enumerate}
\end{Proposition}
\begin{proof}
Note that $f^2\in\Hom_G(V, V)$.
By Schur's lemma, $f^2=k\Id_V$ for some $k\in\mathbb{C}$.
On the other hand,
\begin{gather*}
f(kv)=f(f(f(v)))=kf(v).
\end{gather*}
But $f$ is an anti-linear map on~$V$.
It follows that $k=\overline{k}$ and hence $k\in\mathbb{R}$.
For part (2), we may f\/irst simply pick an isomorphism $f\in\Hom_G(V,
\sigma_G^*\overline{V})$.
Then $f^2=k\Id_V$ for some $k\in\mathbb{R}^\times$.
Schur's lemma implies that any $g\in\Hom_G(V,
\sigma_G^*\overline{V})$ must be of the form $g=cf$ for some $c\in\mathbb{C}$.
Then $g^2=cf\circ cf=c\overline{c}f^2=|c|^2k\Id_V$.
Consequently, if $k$ is positive, we choose $c=\frac{1}{\sqrt{k}}$ so that $g^2=\Id_V$; if $k$ is negative, we choose $c=\frac{1}{\sqrt{-k}}$ so that $g^2=-\Id_V$.
\end{proof}

\begin{Proposition}
\label{realrepclassification}
Let~$V$ be an irreducible Real representation of~$G$.
\begin{enumerate}\itemsep=0pt
\item[$1.$]
The commuting field of~$V$ is isomorphic to $\mathbb{R}$ iff~$V$ is an irreducible complex representation and there
exists $f\in\Hom_G(V,
\sigma_G^*\overline{V})$ such that $f^2=\Id_V$.
\item[$2.$]
The commuting field of~$V$ is isomorphic to $\mathbb{C}$ iff $V\cong W\oplus
\sigma_G^*\overline{W}$ as complex~$G$-representa\-tions,
where $W$ is an irreducible complex~$G$-representation and $W\ncong
\sigma_G^*\overline{W}$, and $\sigma_V(w_1, w_2)$ $=(w_2, w_1)$.
\item[$3.$]
The commuting field of~$V$ is isomorphic to $\mathbb{H}$ iff $V\cong W\oplus
\sigma_G^*\overline{W}$ as complex~$G$-repre\-sen\-ta\-tions,
where $W$ is an irreducible complex~$G$-representation and there exists
$f\in\Hom_G(V,\sigma_G^*\overline{V})$ such that $f^2=-\Id_V$, and $\sigma_V(w_1, w_2)=(w_2, w_1)$.
\end{enumerate}
\end{Proposition}

\begin{proof}
One can easily establish the above proposition
by modifying the proof of Proposition~3 in Appendix~2 of~\cite{B}, which
is a~special case of the above proposition
where $\sigma_G$ is trivial.
\end{proof}
\begin{Proposition}\label{realforgetfulinj}\quad
\begin{enumerate}\itemsep=0pt
\item[$1.$]
The map
$i: RR(G)\to R(G)$ which forgets the Real structure is injective.
\item[$2.$]
Any complex~$G$-representation~$V$ which is a~Real representation can only possess a~unique Real structure up to
isomorphisms of Real~$G$-representations.
\end{enumerate}
\end{Proposition}
\begin{proof}
Let $\rho: R(G)\to RR(G)$ be the map
\begin{gather*}
[V]\mapsto \big[\big(V\oplus\iota_G^*\overline{V},\sigma_{V\oplus\iota_G^*\overline{V}}\big)\big],
\end{gather*}
where $\sigma_{V\oplus\iota_G^*\overline{V}}(u, w)=(w, u)$.
Let $[(V,
\sigma_V)]\in RR(G)$.
Then
\begin{gather*}
\rho\circ i([(V,\sigma_V)])=\big[\big(V\oplus\iota_G^*\overline{V},\sigma_{V\oplus\iota_G^*\overline{V}}\big)\big].
\end{gather*}
We claim that $[(V\oplus V,\sigma_V\oplus\sigma_V)]=[(V\oplus\iota_G^*\overline{V},
\sigma_{V\oplus\iota_G^*\overline{V}})]$, which is easily seen to be true because of the Real~$G$-representation isomorphism
\begin{gather*}
\begin{split}
& f: \ V\oplus \iota_G^*\overline{V}  \to  V\oplus V,
\\
& \phantom{f:{}} \ (u, w)  \mapsto  (u+\sigma_V(w), i(u-\sigma_V(w))).
\end{split}
\end{gather*}
It follows that $\rho\circ i$ amounts to multiplication by 2 on $RR(G)$, and is therefore injective because $RR(G)$ is
a~free abelian group generated by irreducible Real representations.
Hence $i$ is injective.
(2) is simply a~restatement of (1).
\end{proof}
Proposition~\ref{realforgetfulinj} makes it legitimate to regard $RR(G)$ as a~subring of $R(G)$.
From now on we view $[V]\in R(G)$ as an element in $RR(G)$ if~$V$ possesses a~compatible Real structure.
\begin{Proposition}
\label{formclassification}
Let~$G$ be a~compact Real Lie group.
Let~$V$ be an irreducible complex representation of~$G$.
Then
\begin{enumerate}\itemsep=0pt
\item[$1.$]
$[V]\in RR(G, \mathbb{R})$ iff there exists a~$G$-invariant symmetric nondegenerate bilinear form $B: V\times
\sigma_G^*V\to\mathbb{C}$.
\item[$2.$]
$[V\oplus V]\in RR(G, \mathbb{H})$ iff there exists a~$G$-invariant skew-symmetric nondegenerate bilinear form $B:
V\times
\sigma_G^*V\to\mathbb{C}$.
\item[$3.$]
$[V\oplus
\sigma_G^*\overline{V}
]\in RR(G, \mathbb{C})$ iff there does not exist any~$G$-invariant nondegenerate bilinear form on $V\times
\sigma_G^*V$.
\end{enumerate}
\end{Proposition}
\begin{proof}
By Proposition~\ref{realrepclassification},~$V$ is a~Real representation of real type if\/f there exists
$f\in\Hom(V,
\sigma_G^*\overline{V})^G$ such that $f^2=\Id_V$.
One can def\/ine a~bilinear form $B: V\times
\sigma_G^*V\to\mathbb{C}$ by
\begin{gather}
\label{bilinearformtype}
B(v_1, v_2)=\langle v_1, f(v_2)\rangle,
\end{gather}
where $\langle \,, \,\rangle$ is a~$G$-invariant Hermitian inner product on~$V$.
It can be easily seen that $B$ is~$G$-invariant, symmetric and non-degenerate.
Conversely, given a~$G$-invariant symmetric non-degenerate bilinear form on $V\times
\sigma_G^*V$ and using equation~\eqref{bilinearformtype}, we can def\/ine $f\in\Hom(V,
\sigma_G^*\overline{V})^G$, which squares to identity.
Part (2) and (3) follow similarly.
\end{proof}
Proposition~\ref{formclassification} leads to the following

\begin{Definition}
\label{cplxrepclassification}
Let~$G$ be a~compact Real Lie group.
Let~$V$ be an irreducible complex representation of~$G$.
Def\/ine, with respect to $\sigma_G$,
\begin{enumerate}\itemsep=0pt
\item
$V$ to be of real type if there exists a~$G$-invariant symmetric nondegenerate bilinear form $B: V\times
\sigma_G^*V\to\mathbb{C}$.
\item
$V$ to be of quaternionic type if there exists a~$G$-invariant skew-symmetric nondegenerate bilinear form $B: V\times
\sigma_G^*V\to\mathbb{C}$.
\item
$V$ to be of complex type if $V\ncong
\sigma_G^*\overline{V}$.
\end{enumerate}
The abelian group generated by classes of irreducible complex representation of type $\mathbb{F}$ is denoted by $R(G,
\mathbb{F})$.
\end{Definition}

\begin{Definition}
If~$V$ is a~complex~$G$-representation equipped with an anti-linear endomor\-phism~$J_V$ such that $J_V(gv)=
\sigma_G(g)J(v)$ and $J^2=-\Id
_V$, then we say~$V$ is a~\emph{Quaternionic representation} of~$G$.
Let $\mathcal{R}\text{ep}_\mathbb{H}(G)$ be the category of Quaternionic representations of~$G$.
A~morphism between~$V$ and $W\in \mathcal{R}\text{ep}_\mathbb{H}(G)$ is a~linear transformation from~$V$ to $W$ which commutes
with~$G$ and respect $J_V$ and $J_W$.
We denote $\Mor(V, W)$ by $\Hom_G(V, W)^{(J_V, J_W)}$.
An irreducible Quaternionic representation is an irreducible object in $\mathcal{R}\text{ep}_\mathbb{H}(G,
\sigma_G)$.
The Quaternionic representation group of~$G$, denoted by $RH(G)$, is the Grothendieck group of $\mathcal{R}\text{ep}_\mathbb{H}(G)$.
Let $RH(G, \mathbb{F})$ be the abelian group generated by the isomorphism classes of irreducible Quaternionic
representations with $\mathbb{F}$ as the commuting f\/ield.
\end{Definition}
\begin{Remark}
The tensor product of two Quaternionic representations~$V$ and $W$ is a~Real representation as $J_V\otimes J_W$ is an
anti-linear involution which is compatible with $\sigma_G$.
Similarly the tensor product of a~Real representation and a~Quaternionic representation is a~Quaternionic
representation.
To put it succinctly, $RR(G)\oplus RH(G)$ is a~$\mathbb{Z}
_2$-graded ring, with $RR(G)$ being the degree $0$ piece and $RH(G)$ the degree $-1$ piece.
Later on we will assign $RH(G)$ with a~dif\/ferent degree so as to be compatible with the description of the coef\/f\/icient
ring $KR^*_G(\pt)$.
\end{Remark}
\begin{Proposition}
\label{quatrepclassification}
$RH(G)$, as an abelian group, is generated by the following
\begin{enumerate}\itemsep=0pt
\item[$1.$]
$[V\oplus
\sigma_G^*\overline{V}]$ where~$V$ is an irreducible complex representation of real type with $J(u, w)=(-w, u)$.
Its commuting field is $\mathbb{H}$.
\item[$2.$]
$[V\oplus\sigma_G^*\overline{V}]$, where~$V$ is an irreducible complex representation of complex type with $J(u, w)=(-w, u)$.
Its commuting field is $\mathbb{C}$.
\item[$3.$]
$[V]$, where~$V$ is an irreducible complex representation of quaternionic type.
Its commuting field is $\mathbb{R}$.
\end{enumerate}
\end{Proposition}
\begin{proof}
The proof proceeds in the same fashion as does the proof for Proposition~\ref{realrepclassification}.
\end{proof}
\begin{Corollary}\label{3by3table}\quad
\begin{enumerate}\itemsep=0pt
\item[$1.$]
We have that $RR(G, \mathbb{R})\cong RH(G, \mathbb{H})\cong R(G, \mathbb{R})$, $RH(G, \mathbb{R})\cong RR(G,\mathbb{H})$
$\cong R(G, \mathbb{H})$ and $RR(G, \mathbb{C})\cong RH(G, \mathbb{C})$ as abelian groups.
\item[$2.$]
$RR(G)$ is isomorphic to $RH(G)$ as abelian groups.
\end{enumerate}
\end{Corollary}

\begin{proof}
The result follows easily from Proposition~\ref{realrepclassification}, Def\/inition~\ref{cplxrepclassification} and
Proposition~\ref{quatrepclassification}.
\end{proof}
\begin{Proposition}\label{quatforgetfulinj}\quad
\begin{enumerate}\itemsep=0pt
\item[$1.$]
The map $j: RH(G)\to R(G)$ which forgets the Quaternionic structure is injective.
\item[$2.$]
Any complex~$G$-representation which is a~Quaternionic representation can only possess a~unique Quaternionic structure
up to isomorphisms of Quaternionic~$G$-representations.
\end{enumerate}
\end{Proposition}
\begin{proof}
The proof proceeds in the same fashion as does the proof for Proposition~\ref{realforgetfulinj}.
It suf\/f\/ices to show that, if $\eta: R(G)\to RH(G)$ is the map $[V]\mapsto [(V\oplus
\sigma_G^*\overline{V},
\sigma_{V\oplus\iota_G^*\overline{V}})]$, then $\eta\circ j$ amounts to multiplication by 2, i.e.\
$(V\oplus\sigma_G^*\overline{V},\sigma_{V\oplus\sigma_G^*\overline{V}})\cong(V\oplus V, J\oplus J)$, where $\sigma_{V\oplus\sigma_G^*\overline{V}}(u, w)=(-w, u)$.
This is true because of the Quaternionic~$G$-representation isomorphism
\begin{gather*}
f: \  V\oplus\iota_G^*\overline{V}  \to   V\oplus V,\\
\hphantom{f: {}}  \ (u, w) \mapsto  (u+Jw, i(u-Jw)). \tag*{\qed}
\end{gather*}
\renewcommand{\qed}{}
\end{proof}

\begin{Example}
We shall illustrate the similarities and dif\/ferences of the various aforementioned representation groups with an
example.
Let $G=Q_8\times C_3$, the direct product of the quaternion group and the cyclic group of order 3, equipped with the
trivial involution.
There are 5 irreducible complex representations of $Q_8$, namely, the 4 1-dimensional representations which become
trivial on restriction to the center $Z$ of $Q_8$ and descend to the 4 1-dimensional representations of $Q_8/Z\cong
\mathbb{Z}_2\oplus\mathbb{Z}_2$, and the 2-dimensional faithful representation.
We denote these representations by $1_{Q_8}$, $\rho_{(0, 1)}$, $\rho_{(1, 0)}$, $\rho_{(1, 1)}$ and~$\rho_Q$ respectively.
Similarly, we let~$1_{C_3}$,~$\rho_\zeta$ and~$\rho_{\zeta^2}$ be the three 1-dimensional complex irreducible
representations of~$C_3$.
It can be easily seen that
\begin{gather*}
R(Q_8, \mathbb{R})=\mathbb{Z}\cdot[1_{Q_8}]\oplus\mathbb{Z}\cdot[\rho_{(0, 1)}]\oplus\mathbb{Z}\cdot[\rho_{(1,
0)}]\oplus\mathbb{Z}\cdot[\rho_{(1, 1)}],
\\
R(Q_8, \mathbb{C})=0,
\\
R(Q_8, \mathbb{H})=\mathbb{Z}\cdot[\rho_Q],
\\
R(C_3, \mathbb{R})=\mathbb{Z}\cdot[1_{C_3}],
\\
R(C_3, \mathbb{C})=\mathbb{Z}\cdot[\rho_\zeta]\oplus\mathbb{Z}\cdot[\rho_{\zeta^2}],
\\
R(C_3, \mathbb{H})=0.
\end{gather*}
It follows that
\begin{gather*}
R(G, \mathbb{R})=\bigoplus_{x\in \{1_{Q_8}, \rho_{(0, 1)}, \rho_{(1, 0)}, \rho_{(1,
1)}\}}\mathbb{Z}\cdot[x\widehat{\otimes}1_{C_3}]\cong\mathbb{Z}^4,
\\
R(G, \mathbb{C})=\bigoplus_{
\begin{subarray}
{c}{x\in\{1_{Q_8}, \rho_{(0, 1)}, \rho_{(1, 0)}, \rho_{(1, 1)}, \rho_Q\}}
\\
{y\in\{\rho_\zeta, \rho_{\zeta^2}\}}
\end{subarray}}\mathbb{Z}\cdot[x\widehat{\otimes}y]\cong\mathbb{Z}^{10},
\\
R(G, \mathbb{H})=\mathbb{Z}\cdot[\rho_Q\widehat{\otimes}1_{C_3}]\cong\mathbb{Z},
\\
RR(G, \mathbb{R})=\bigoplus_{x\in \{1_{Q_8}, \rho_{(0, 1)}, \rho_{(1, 0)}, \rho_{(1,
1)}\}}\mathbb{Z}\cdot[x\widehat{\otimes}1_{C_3}]\cong\mathbb{Z}^4,
\\
RR(G, \mathbb{C})=\bigoplus_{x\in \{1_{Q_8}, \rho_{(0, 1)}, \rho_{(1, 0)}, \rho_{(1, 1)},
\rho_Q\}}\mathbb{Z}\cdot[x\widehat{\otimes}\rho_\zeta\oplus\overline{x}\widehat{\otimes}\rho_{\zeta^2}]\cong\mathbb{Z}^5,
\\
RR(G, \mathbb{H})=\mathbb{Z}\cdot[\rho_Q\widehat{\otimes}1_{C_3}\oplus
\overline{\rho_Q}\widehat{\otimes}1_{C_3}]\cong\mathbb{Z},
\\
RH(G, \mathbb{R})=\mathbb{Z}\cdot[\rho_Q\widehat{\otimes}1_{C_3}]\cong\mathbb{Z},
\\
RH(G, \mathbb{C})=\bigoplus_{x\in \{1_{Q_8}, \rho_{(0, 1)}, \rho_{(1, 0)}, \rho_{(1, 1)},
\rho_Q\}}\mathbb{Z}\cdot[x\widehat{\otimes}\rho_\zeta\oplus\overline{x}\widehat{\otimes}\rho_{\zeta^2}]\cong\mathbb{Z}^5,
\\
RH(G, \mathbb{H})=\bigoplus_{x\in \{1_{Q_8}, \rho_{(0, 1)}, \rho_{(1, 0)}, \rho_{(1,
1)}\}}\mathbb{Z}\cdot[x\widehat{\otimes}1_{C_3}\oplus\overline{x}\widehat{\otimes}1_{C_3}]\cong\mathbb{Z}^4.
\end{gather*}
Some representations above should be equipped with suitable Real or Quaternionic structures given in
Propositions~\ref{realrepclassification} and~\ref{quatrepclassification}.
For example, the Real structure of $\rho_Q\widehat{\otimes}1_{C_3}\oplus\overline{\rho_Q}\widehat{\otimes}1_{C_3}$ in
$RR(G, \mathbb{H})$ is given by swapping the two coordinates.
\end{Example}

\subsection[$KR$-theory]{$\boldsymbol{KR}$-theory}

$KR$-theory was f\/irst introduced by Atiyah in~\cite{At3} and used to derive the 8-periodicity of $KO$-theory from the
2-periodicity of complex $K$-theory.
$KR$-theory was motivated by the index theory of real elliptic operators.
\begin{Definition}\label{krtheory}\quad
\begin{enumerate}\itemsep=0pt
\item
A \emph{Real space} is a~pair $(X,
\sigma_X)$ where $X$ is a~topological space equipped with an involutive homeomorphism $\sigma_X$, i.e.\
$\sigma_X^2=\Id
_X$.
We will sometimes suppress the notation $\sigma_X$ and simply use $X$ to denote the Real space,
if there is no danger of confusion about the involutive
homeomorphism.
A \emph{Real pair}
is a~pair $(X, Y)$ where~$Y$ is a~closed subspace of~$X$ invariant under~$\sigma_X$.
\item
Let $\mathbb{R}^{p, q}$ be the Euclidean space $\mathbb{R}^{p+q}$
equipped with the involution which is identity on the f\/irst~$q$
coordinates and negation on the last $p$-coordinates.
Let $B^{p, q}$ and $S^{p, q}$ be the unit ball and sphere in $\mathbb{R}^{p, q}$ with the inherited involution.
\item
A~\emph{Real vector bundle}
(to be distinguished from the usual real vector bundle) over $X$ is a~complex vector bundle
$E$ over $X$ which itself is also a~Real space with involutive homeomorphism $\sigma_E$ satisfying
\begin{enumerate}\itemsep=0pt
\item
$\sigma_X\circ p=p\circ\sigma_E$, where $p: E\to X$ is the projection map,
\item
$\sigma_E$ maps $E_x$ to $E_{\sigma_X(x)}$ anti-linearly.
\end{enumerate}
A \emph{Quaternionic vector bundle}
(to be distinguished from the usual quaternionic vector bundle) over $X$ is a~complex
vector bundle $E$ over $X$ equipped with an anti-linear lift $\sigma_E$ of $\sigma_X$ such that $\sigma_E^2=-\Id
_E$.
\item
Let $X$ be a~Real space.
The ring $KR(X)$ is the Grothendieck group of the isomorphism classes of Real vector bundles over $X$, equipped with the
usual product structure induced by tensor product of vector bundles over $\mathbb{C}$.
The relative $KR$-theory for a~Real pair $KR(X, Y)$ can be similarly def\/ined.
In general, the graded $KR$-theory ring of the Real pair $(X, Y)$ is given by
\begin{gather*}
KR^*(X, Y):=\bigoplus_{q=0}^7 KR^{-q}(X, Y),
\end{gather*}
where
\begin{gather*}
KR^{-q}(X, Y):=KR\big(X\times B^{0,q}, X\times S^{0, q}\cup Y\times B^{0,q}\big).
\end{gather*}
The ring structure of $KR^*$ is extended from that of~$KR$, in a~way analogous to the case of complex $K$-theory.
The number of graded pieces, which is~8, is a~result of Bott periodicity for $KR$-theory (cf.~\cite{At3}).
\end{enumerate}
\end{Definition}

Note that when $\sigma_X=\Id_X$, then $KR(X)\cong KO(X)$.
On the other hand, if $X\times \mathbb{Z}_2$ is given the involution which swaps the two copies of $X$, then
$KR(X\times\mathbb{Z}_2)\cong K(X)$.
Also, if $X$ is equipped with the trivial involution, then $KR(X\times S^{2,0})\cong KSC(X)$, the Grothendieck group of
homotopy classes of self-conjugate bundles over $X$ (cf.~\cite{At3}).
In this way, it is natural to view $KR$-theory as a~unifying thread of $KO$-theory, $K$-theory and $KSC$-theory.

On top of the Real structure, we may further add compatible group actions and def\/ine equivariant $KR$-theory.
\begin{Definition}\label{eqkrtheory}\quad
\begin{enumerate}\itemsep=0pt
\item
A \emph{Real}~$G$-\emph{space} $X$ is a~quadruple $(X, G,
\sigma_X, \sigma_G)$ where a~group~$G$ acts on~$X$ and~$\sigma_G$ is an involutive automorphism of~$G$ such that
\begin{gather*}
\sigma_X(g\cdot x)=\sigma_G(g)\cdot\sigma_X(x).
\end{gather*}
\item
A \emph{Real}~$G$-\emph{vector bundle} $E$ over a~Real~$G$-space $X$ is a~Real vector bundle and a~$G$-bundle over $X$,
and it is also a~Real~$G$-space.
\item
In a~similar spirit, one can def\/ine equivariant $KR$-theory $KR_G^*(X, Y)$.
Notice that the~$G$-actions on $B^{0,q}$ and $S^{0,q}$ in the def\/inition of $KR_G^{-q}(X, Y)$ are trivial.
\end{enumerate}
\end{Definition}

\begin{Definition}\quad
\begin{enumerate}\itemsep=0pt
\item
Let $K^*(+)$ be the complex $K$-theory of a~point extended to a~$\mathbb{Z}_8$-graded algebra over
$K^0(\pt)\cong\mathbb{Z}$, i.e.\
$\displaystyle K^*(+)\cong\mathbb{Z}[\beta]\left/\beta^4-1\right.$.
Here $\beta\in K^{-2}(+)$ is the class of the reduced canonical bundle on $\mathbb{CP}^1\cong S^2$.
\item
Let $\overline{\sigma_X^*}$ be the map def\/ined on (equivariant) vector bundles on $X$ by $\overline{\sigma_X^*}
E:=
\sigma_X^*\overline{E}$.
The involution induced by $\overline{\sigma_X^*}$ on $K^*_G(X)$ is also denoted by $\overline{\sigma_X^*}$ for simplicity.
\end{enumerate}
\end{Definition}
In the following proposition, we collect, for reader's convenience, some basic results of $KR$-theory
(cf.~\cite[Section~2]{Se}), some of which are stated in the more general context of equivariant $KR$-theory.
\begin{Proposition}\label{krprelim}\quad
\begin{enumerate}\itemsep=0pt
\item[$1.$]
We have
\begin{gather*}
KR^*(\pt)\cong\mathbb{Z}[\eta, \mu]\big/\big(2\eta, \eta^3, \mu\eta, \mu^2-4\big) ,
\end{gather*}
where $\eta\in KR^{-1}(\pt)$, $\mu\in KR^{-4}(\pt)$ represents the reduced Hopf bundles of
$\mathbb{R}\mathbb{P}^1$ and $\mathbb{H}\mathbb{P}^1$ respectively.
\item[$2.$]
Let $c: KR^*_G(X)\to K^*_G(X)$ be the homomorphism which forgets the Real structure of Real vector bundles, and $r:
K^*_G(X)\to KR^*_G(X)$ be the realification map defined by $[E]\mapsto [E\oplus
\sigma_G^*\sigma_X^*\overline{E}
]$.
Then we have the following relations
\begin{enumerate}\itemsep=0pt
\item[$(a)$]
$c(1)=1$, $c(\eta)=0$, $c(\mu)=2\beta^2$, where $\beta\in K^{-2}(\pt)$ is the Bott class,
\item[$(b)$]
$r(1)=2$, $r(\beta)=\eta^2$, $r(\beta^2)=\mu$, $r(\beta^3)=0$,
\item[$(c)$]
$r(xc(y))=r(x)y$, $cr(x)=x+
\sigma_G^*\overline{\sigma_X^*}
x$ and $rc(y)=2y$ for $x\in K^*_G(X)$ and $y\in KR^*_G(X)$, where $K^*_G(X)$ is extended to a~$\mathbb{Z}_8$-graded
algebra by Bott periodicity.
\end{enumerate}
\end{enumerate}
\end{Proposition}
\begin{proof}
(1) is given in~\cite[Section~2]{Se}.
The proof of (2) is the same as in the nonequivariant case, which is given in~\cite{At3}.
\end{proof}

\begin{Definition}
A Quaternionic~$G$-vector bundle over a~Real space $X$ is a~complex vector bundle $E$ equipped with an anti-linear
vector bundle endomorphism $J$ on $E$ such that $J^2=-\Id_E$ and $J(g\cdot v)=
\sigma_G(g)\cdot J(v)$.
Let $KH^*_G(X)$ be the corresponding $K$-theory constructed using Quaternionic~$G$-bundles over $X$.
\end{Definition}
By generalizing the discussion preceding Lemma 5.2 in~\cite{Se} to the equivariant and graded setting, we def\/ine
a~natural transformation
\begin{gather*}
t: \ KH_G^{-q}(X)\to KR^{-q-4}_G(X)
\end{gather*}
which sends
\begin{gather*}
0\longrightarrow E_1
\stackrel{f}
{\longrightarrow} E_2\longrightarrow 0
\end{gather*}
to
\begin{gather*}
0\longrightarrow \pi^*(\mathbb{H}\otimes_\mathbb{C} E_1)\stackrel{g}
{\longrightarrow}\pi^*(\mathbb{H}\otimes_\mathbb{C}E_2)\longrightarrow 0,
\end{gather*}
where
\begin{enumerate}\itemsep=0pt
\item[1)]
$E_i$, $i=1, 2$ are equivariant Quaternionic vector bundles on $X\times\mathbb{R}^{0, q}$ equipped with the Quaternionic
structures $J_{E_i}$,
\item[2)]
$f$ is an equivariant Quaternionic vector bundle homomorphism which is an isomorphism outside $X\times\{0\}$,
\item[3)]
$\pi: X\times\mathbb{R}^{0, q+4}\to X\times\mathbb{R}^{0, q}$ is the projection map,
\item[4)]
$\mathbb{H}\otimes_\mathbb{C}E_i$ is the equivariant Real vector bundles equipped with the Real structure $J\otimes J_{E_i}$,
\item[5)]
$g$ is an equivariant Real vector bundle homomorphism def\/ined by $g(v, w\otimes e)=(v, vw\otimes f(e))$.
\end{enumerate}

One can easily show by generalizing the discussion in the last section of~\cite{AS} that
\begin{Proposition}
\label{shiftbyfour}
$t$ is an isomorphism.
\end{Proposition}

\subsection[The module structure of $KR$-theory of compact simply-connected Lie groups]{The
module structure of $\boldsymbol{KR}$-theory\\ of compact simply-connected Lie groups}
The following structure theorem for $KR$-theory, due to Seymour, is crucial in his computation of
$KR^*(\pt)$-module structure of $KR^*(G)$.
\begin{Theorem}[{\cite[Theorem 4.2]{Se}}]\label{strthm}
Suppose that $K^*(X)$ is a~free abelian group and decomposed by the involution $\overline{\sigma_X^*}$ into the following summands
\begin{gather*}
K^*(X)=M_+\oplus M_-\oplus T\oplus \overline{\sigma_X^*}T,
\end{gather*}
where $\overline{\sigma_X^*}$ is identity on $M_+$ and negation on $M_-$.
Suppose further that there exist $h_1, \dots, h_n\in KR^*(X)$ such that $c(h_1), \dots, c(h_n)$ form a~basis of the
$K^*(+)$-module $K^*(+)\otimes(M_+\oplus M_-)$.
Then, as $KR^*(\pt)$-modules,
\begin{gather*}
KR^*(X)\cong F\oplus r(K^*(+)\otimes T),
\end{gather*}
where $F$ is the free $KR^*(\pt)$-module generated by $h_1, \dots, h_n$.
\end{Theorem}
\begin{Remark}
If $T=0$, then the conditions in Theorem~\ref{strthm} are equivalent to $K^*(X)$ being free abelian and $c: KR^*(X)\to
K^*(X)$ being surjective.
In this special case the theorem
implies that the map $KR^*(X)\otimes_{KR^*(\pt)}K^*(\pt)\to K^*(X)$ def\/ined
by $a\otimes b\mapsto c(a)\cdot b$ is a~ring isomorphism.
This smacks of the def\/inition of weakly equivariant formality (cf.\ Remark~\ref{weakequivformality}) and inspires us to
def\/ine a~similar notion for equivariant $KR$-theory (cf.\ Def\/inition~\ref{realequivformal}).
We say a~real space is \emph{real formal} if it satisf\/ies the conditions of Theorem~\ref{strthm}.
\end{Remark}
\begin{Definition}
Let $\sigma_\mathbb{R}$ be the complex conjugation of $U(n)$ or $U(\infty)$,
and $\sigma_\mathbb{H}$ be the symplectic type involution $g\mapsto J_m\overline{g}J_m^{-1}$ on $U(2m)$, or $U(2\infty)$.
\end{Definition}
For any Real space $X$, $KR^{-1}(X)$ is isomorphic to the abelian group of equivariant homotopy classes of maps from $X$
to $U(\infty)$ which respect $\sigma_X$ and $\sigma_\mathbb{R}$ on $U(\infty)$.
Similarly, $KR^{-5}(X)$, which is isomorphic to $KH^{-1}(X)$ by Proposition~\ref{shiftbyfour}, is isomorphic to the
abelian group of equivariant homotopy classes of maps from $X$ to $U(2\infty)$
which respect $\sigma_X$ and $\sigma_\mathbb{H}$ on $U(2\infty)$
(cf.\ remarks in the last two paragraphs of Appendix of~\cite{Se}).
We can def\/ine maps analogous to those in Def\/inition~\ref{betamap} in the context of $KR$-theory.
\begin{Definition}
%\label{rbetamap}
Let $\delta_\mathbb{R}: RR(G)\to KR^{-1}(G)$ and $\delta_\mathbb{H}: RH(G)\to KR^{-5}(G)$ be group homomorphisms which
send a~Real (resp.\
Quaternionic) representation to the $KR$-theory element represented by its homotopy class.
\end{Definition}
\begin{Proposition}
If $\rho\in RR(G)$, then $\delta_\mathbb{R}(\rho)$ is represented by the complex of vector bundles in
Proposition~{\rm \ref{-1_class}} equipped with the Real structure given by
\begin{gather*}
\iota: \ G\times\mathbb{R}\times V  \to  G\times\mathbb{R}\times V,
\\
\phantom{\iota: {}} \
(g, t, v)  \mapsto  (\sigma_G(g), t, \overline{v}).
\end{gather*}
If $\rho\in RH(G)$, then $\delta_\mathbb{H}(\rho)$ can be similarly represented, with the Real structure replaced by the
Quaternionic structure.
\end{Proposition}

\emph{From this point on until the end of this section, we further assume that~$G$ is connected and simply-connected unless
otherwise specified}.
It is known that $R(G)$ is a~polynomial ring over $\mathbb{Z}$ generated by fundamental representations, which are
permuted by $\overline{\sigma_{G}^*}$ (cf.\ \cite[Lemma~5.5]{Se}).
Let
\begin{gather*}
R(G)\cong\mathbb{Z}\big[\varphi_1, \dots, \varphi_r, \theta_1, \dots, \theta_s, \gamma_1, \dots, \gamma_t, \overline{\sigma_G^*}
\gamma_1, \dots, \overline{\sigma_G^*}\gamma_t\big],
\end{gather*}
where $\varphi_i\in RR(G, \mathbb{R})$, $\theta_j\in RH(G, \mathbb{R})$, $\gamma_k\in R(G, \mathbb{C})$.
Then $K^*(G)$, as a~free abelian group, is generated by square-free monomials
in $\delta(\varphi_1), {\dots},\delta(\varphi_r)$,
$\delta(\theta_1), {\dots}, \delta(\theta_s)$,
$\delta(\gamma_1), {\dots}, \delta(\gamma_t)$,
$\delta(\overline{\sigma_G^*}\gamma_1)$, ${\dots}, \delta(\overline{\sigma_G^*}\gamma_t)$.
Using Theorem~\ref{strthm}, Seymour obtained

\begin{Theorem}[{\cite[Theorem 5.6]{Se}}]
\label{krtheorymodofg}\quad
\begin{enumerate}\itemsep=0pt
\item[$1.$]
Suppose that $\overline{\sigma_G^*}$ acts as identity on $R(G)$, i.e.\
any irreducible Real representation of~$G$ is either of real type or quaternionic type.
Then as $KR^*(\pt)$-modules,
\begin{gather*}
KR^*(G)\cong\wedge_{KR^*(\pt)}(\delta_\mathbb{R}(\varphi_1), \dots, \delta_\mathbb{R}(\varphi_r),
\delta_\mathbb{H}(\theta_1), \dots, \delta_\mathbb{H}(\theta_s)).
\end{gather*}
\item[$2.$]
More generally, $c(\delta_\mathbb{R}(\varphi_i))=\delta(\varphi_i)$,
$c(\delta_\mathbb{H}(\theta_j))=\beta^2\cdot\delta(\theta_j)$, and there exist $\lambda_1, \dots, \lambda_t\in KR^0(G)$
such that $c(\lambda_k)=\beta^3\cdot\delta(\gamma_k)\delta(\overline{\sigma_G^*}
\gamma_k)$, and
\begin{gather*}
KR^*(G)\cong P\oplus T\cdot P
\end{gather*}
as $KR^*(\pt)$-module, where
\begin{itemize}\itemsep=0pt
\item
$P\cong\bigwedge_{KR^*(\pt)}(\delta_\mathbb{R}(\varphi_1), \dots, \delta_\mathbb{R}(\varphi_r),
\delta_\mathbb{H}(\theta_1), \dots, \delta_\mathbb{H}(\theta_s), \lambda_1, \dots, \lambda_t)$,
\item
$T$ is the additive abelian group generated by the set
\begin{gather*}
\{r(\beta^i\cdot\delta(\gamma_1)^{\varepsilon_1}\cdots\delta(\gamma_t)^{\varepsilon_t}\delta(\overline{\sigma_G^*}
\gamma_1)^{\nu_1}\cdots\delta(\overline{\sigma_G^*}\gamma_t)^{\nu_t})\},
\end{gather*}
where $\varepsilon_1, \dots, \varepsilon_t$, $\nu_1, \dots, \nu_t$ are either $0$ or $1$, $\varepsilon_k$ and~$\nu_k$ are not
equal to 1 at the same time for $1\leq k\leq t$, and the first index $k_0$ where $\varepsilon_{k_0}=1$ is less than the
first index~$k_1$ where $\nu_{k_1}=1$.
\end{itemize}
Moreover,
\begin{enumerate}\itemsep=0pt
\item[$(a)$]
$\lambda_k^2=0$ for all $1\leq k\leq t$,
\item[$(b)$]
$\delta_\mathbb{R}(\varphi_i)^2$ and $\delta_\mathbb{H}(\theta_j)^2$ are divisible by $\eta$.
\end{enumerate}
\end{enumerate}
\end{Theorem}

\begin{Definition}
Let $\omega_t:=\delta_{\epsilon_t, 1-\nu_t}$ and
\begin{gather*}
r_{i, \varepsilon_1, \dots, \varepsilon_t, \nu_1, \dots,
\nu_t}:=r\left(\beta^i\cdot\delta(\gamma_1)^{\varepsilon_1}\cdots\delta(\gamma_t)^{\varepsilon_t}\delta(\overline{\sigma_G^*}
(\gamma_1))^{\nu_1}\cdots\delta\big(\overline{\sigma_G^*}(\gamma_t)\big)^{\nu_t}\right)\in T.
\end{gather*}
\end{Definition}
\begin{Corollary}\label{krglooseend}\quad
\begin{enumerate}\itemsep=0pt
\item[$1.$]
$KR^*(G)$ is generated by $\delta_\mathbb{R}(\varphi_1),  {\dots}, \delta_\mathbb{R}(\varphi_r)$,
$\delta_\mathbb{H}(\theta_1),  {\dots}, \delta_\mathbb{H}(\theta_s)$, $\lambda_1,  {\dots}, \lambda_t\!$ and $r_{i,
\varepsilon_1, {\dots}, \varepsilon_t, \nu_1, {\dots}, \nu_t}$ $\in T$ as an algebra over $KR^*(\pt)$.
\item[$2.$]
\begin{gather*}
r_{i, \varepsilon_1, \dots, \varepsilon_t, \nu_1, \dots, \nu_t}^2=
\begin{cases}
\eta^2\lambda_1^{\omega_1}\cdots\lambda_t^{\omega_t},
&\text{if }r_{i, \varepsilon_1, \dots, \varepsilon_t, \nu_1, \dots, \nu_t}\text{ is of degree }-1\text{ or }-5,
\\
\pm\mu\lambda_1^{\omega_1}\cdots\lambda_t^{\omega_t},
&\text{if }r_{i, \varepsilon_1, \dots, \varepsilon_t, \nu_1, \dots, \nu_t}\text{ is of degree }-2\text{ or }-6,
\\
0
&\text{otherwise}.
\end{cases}
\end{gather*}
The sign depends on $i$, $\varepsilon_1, \dots, \varepsilon_t$, $\nu_1, \dots, \nu_t$ and can be determined using formulae
from~$(2)$ of Proposition~{\rm \ref{krprelim}}.
\item[$3.$]
$r_{i, \varepsilon_1, \dots, \varepsilon_t, \nu_1, \dots, \nu_t}\eta=0$, and $r_{i, \varepsilon_1, \dots, \varepsilon_t,
\nu_1, \dots, \nu_t}\mu=2r_{i+2, \varepsilon_1, \dots, \varepsilon_t, \nu_1, \dots, \nu_t}$.
\end{enumerate}
\end{Corollary}

\begin{proof}
The Corollary follows easily from the various properties of the realif\/ication map and the complexif\/ication map in
Proposition~\ref{krprelim}, and the fact that $c(\lambda_k)=\beta^3\cdot\delta(\gamma_k)\delta(\overline{\sigma_G^*}
\gamma_k)$.
For example,
\begin{gather*}
r_{i, \varepsilon_1, \dots, \varepsilon_t, \nu_1, \dots,
\nu_t}\eta=r\big(\beta^i\cdot\delta(\gamma_1)^{\varepsilon_1}\cdots\delta(\gamma_t)^{\varepsilon_t}\delta(\overline{\sigma_G^*}
(\gamma_1))^{\nu_1}\cdots\delta\big(\overline{\sigma_G^*}(\gamma_t)\big)^{\nu_t}c(\eta)\big)
=0,
\\
r_{i, \varepsilon_1, \dots, \varepsilon_t, \nu_1, \dots,\nu_t}\mu
=r\big(\beta^i\cdot\delta(\gamma_1)^{\varepsilon_1}\cdots\delta(\gamma_t)^{\varepsilon_t}\delta(\overline{\sigma_G^*}
(\gamma_1))^{\nu_1}\cdots\delta\big(\overline{\sigma_G^*}(\gamma_t)\big)^{\nu_t}c(\mu)\big)
\\
\phantom{r_{i, \varepsilon_1, \dots, \varepsilon_t, \nu_1, \dots,\nu_t}\mu}
=r_{i+2, \varepsilon_1, \dots, \varepsilon_t, \nu_1, \dots, \nu_t}.\tag*{\qed}
\end{gather*}
\renewcommand{\qed}{}
\end{proof}

In fact Theorems~\ref{hodgkin} and~\ref{strthm} also yield the following description of module structure of
$KR$-theory of a~compact connected Real Lie group with torsion-free fundamental group with a~restriction on the types of
the Real representations.
\begin{Theorem}
\label{spkrtheorymodofg}
Let~$G$ be a~compact connected real Lie group with $\pi_1(G)$ torsion-free.
Suppose that $R(G, \mathbb{C})=0$, i.e.\
$\overline{\sigma_G^*}$ acts as identity on $R(G)$.
Then $KR^*(G)$ is isomorphic to $\wedge_{KR^*(\pt)}^*(\text{\rm Im}(\widetilde{\delta}_\mathbb{R}),
\text{\rm Im}(\widetilde{\delta}_\mathbb{H}))$ as $KR^*(\pt)$-modules.
\end{Theorem}

As we see from Theorem~\ref{krtheorymodofg} and Corollary~\ref{krglooseend}, to get a~full description of the ring
structure of $KR^*(G)$, it remains to f\/igure out $\delta_\mathbb{R}(\varphi_i)^2$ and $\delta_\mathbb{H}(\theta_j)^2$.
We will, in the end, obtain formulae for the squares by way of computing the ring structure of $KR_G^*(G)$ and applying
the forgetful map.
In particular, we will show that $\delta_\mathbb{R}(\varphi_i)^2$ and $\delta_\mathbb{H}(\theta_j)^2$ in general are
non-zero.
So, unlike the complex $K$-theory, $KR^*(G)$ is not an exterior algebra in general.
Nevertheless, $KR^*(G)$ is not far from being an exterior algebra, in the sense of the following
\begin{Corollary}\quad
\begin{enumerate}\itemsep=0pt
\item[$1.$]
$KR^*(\pt)_2$, which is the ring obtained by inverting the prime~$2$ in $KR^*(\pt)$, is isomorphic to
$\mathbb{Z}\left[\frac{1}{2}, \mu\right]/(\mu^2-4)\cong\mathbb{Z}\left[\frac{1}{2}, \beta^2\right]/((\beta^2)^2-1)$.
\item[$2.$]
Suppose that $R(G, \mathbb{C})=0$.
$KR^*(G)_2$, which is the ring obtained by inverting the prime~$2$ in $KR^*(G)$, is isomorphic to, as
$KR^*(\pt)_2$-algebra
\begin{gather*}
\bigwedge\nolimits_{KR^*(\pt)_2}\big(\delta_\mathbb{R}(\varphi_1), \dots, \delta_\mathbb{R}(\varphi_r),
\delta_\mathbb{H}(\theta_1), \dots, \delta_\mathbb{H}(\theta_s)\big).
\end{gather*}
\end{enumerate}
\end{Corollary}

\section[The coef\/f\/icient ring $KR_G^*(\pt)$]{The coef\/f\/icient ring $\boldsymbol{KR_G^*(\pt)}$}\label{Section3}

In this section, we assume that~$G$ is a~compact Real Lie group, and will prove a~result on the coef\/f\/icient ring
$KR_G^*(\pt)$.
In~\cite{AS}, all graded pieces of $KR_G^*(\pt)$ were worked out using Real Clif\/ford~$G$-modules.
We record them in the following
\begin{Proposition}
\label{ASKRcomp}
$KR^{-q}_G(\pt)$, as abelian groups, for $0\leq q\leq 7$, are isomorphic to $RR(G)$, $RR(G)/\rho(R(G))$,
$R(G)/j(RH(G))$, $0$, $RH(G)$, $RH(G)/\eta(R(G))$, $R(G)/i(RR(G))$ and $0$ respectively, where the maps $i$, $j$, $\rho$, $\eta$
are as in Propositions~{\rm \ref{realforgetfulinj}} and~{\rm \ref{quatforgetfulinj}}.
\end{Proposition}
\begin{Remark}
Note from the above proposition
that $KR_G^0(\pt)\oplus KR^{-4}_G(\pt)\cong RR(G)\oplus RH(G)$.
In this way we can view $RR(G)\oplus RH(G)$ as a~graded ring where $RR(G)$ is of degree 0 and $RH(G)$ of degree $-4$.
\end{Remark}
\begin{Proposition}\label{KR_G}\quad
\begin{enumerate}\itemsep=0pt
\item[$1.$]
Suppose $R(G, \mathbb{C})=0$.
Then the map
\begin{gather*}
f:  \ (RR(G, \mathbb{R})\oplus RH(G, \mathbb{R}))\otimes KR^*(\pt)  \to  KR^*_G(\pt),
\\
\phantom{f: {}} \
\rho_1\otimes x_1\oplus \rho_2\otimes x_2  \mapsto  \rho_1\cdot x_1+\rho_2\cdot x_2
\end{gather*}
is an isomorphism of graded rings.
\item[$2.$]
In general,
\begin{gather*}
f: \ (RR(G, \mathbb{R})\oplus RH(G, \mathbb{R}))\otimes KR^*(\pt)\oplus r(R(G, \mathbb{C})\otimes K^*(+))\to
KR^*_G(\pt),
\\
\phantom{f: {}} \
\rho_1\otimes x_1\oplus \rho_2\otimes x_2\oplus r(\rho_3\otimes \beta^i)\mapsto \rho_1\cdot x_1+\rho_2\cdot
x_2+r(\rho_3\cdot\beta^i)
\end{gather*}
is an isomorphism of graded abelian groups.
\item[$3.$]
If $\rho$ is an irreducible complex representation of complex type, then $\eta r(\beta^i\cdot\rho)=0$ and $\mu
r(\beta^i\cdot\rho)=2r(\beta^{i+2}\cdot\rho)$.
\end{enumerate}
\end{Proposition}

\begin{proof}
The proposition
follows by verifying the isomorphism in dif\/ferent degree pieces against the description in
Proposition~\ref{ASKRcomp}.
For example, in degree 0,
\begin{gather*}
RR(G, \mathbb{R})\otimes KR^0(\pt)\oplus RH(G, \mathbb{R})\otimes KR^{-4}(\pt)\oplus r(R(G,
\mathbb{C})\otimes K^0(+))
\\
\qquad=RR(G, \mathbb{R})\oplus RH(G, \mathbb{R})\otimes \mathbb{Z}\mu\oplus RR(G, \mathbb{C})
\\
\qquad
\cong RR(G, \mathbb{R})\oplus RR(G, \mathbb{H})\oplus RR(G, \mathbb{C})
\\
(\text{if} \ \ [V]\in RH(G, \mathbb{R}),\ \ \text{then} \ \ [V]\cdot\mu=[V\oplus V]\in RR(G, \mathbb{H}))
\\
\qquad=RR(G)
=KR^0_G(\pt).
\end{gather*}
(3) follows from Proposition~\ref{krprelim}.
\end{proof}
\begin{Remark}
In~\cite{AS}, $KR_G^*(X)$, where the~$G$-action is trivial, is given as the following direct sum of abelian groups
\begin{gather*}
RR(G, \mathbb{R})\otimes KR^*(X)\oplus RR(G, \mathbb{C})\otimes KC^*(X)\oplus RR(G, \mathbb{H})\otimes KH^*(X),
\end{gather*}
where $KC^*(X)$ and $KH^*(X)$ are Grothendieck groups of the so-called `Complex vector bundles' and `Quaternionic vector
bundles' of $X$.
We f\/ind Proposition~\ref{KR_G}, which is motivated by this description, better because the ring structure of the
coef\/f\/icient ring is more apparent when cast in this light.
The proposition
is, as we will see in the next section, a~consequence of a~structure theorem of equivariant $KR$-theory
(Theorem~\ref{equivstrthm}), and therefore still holds true if the point is replaced by any general space $X$ with
trivial~$G$-action.
\end{Remark}

\section[Equivariant $KR$-theory rings of compact simply-connected Lie groups]{Equivariant
$\boldsymbol{KR}$-theory rings \\ of compact simply-connected Lie groups}\label{Section4}

Throughout this section we assume that~$G$ is a~compact, connected and simply-connected Real Lie group unless otherwise
specif\/ied.
We will prove the main result of this paper, Theorem~\ref{mainthm2}, which gives the ring structure of $KR_G^*(G)$.
Our strategy is outlined as follows.
\begin{enumerate}\itemsep=0pt
\item
We obtain a~result on the structure of $KR^*_G(G)$ (Corollary~\ref{equivstrthmg}) which is analogous to
Theorem~\ref{krtheorymodofg} and Proposition~\ref{KR_G}.
We def\/ine $\delta_\mathbb{R}^G(\varphi_i)$, $\delta_\mathbb{H}^G(\theta_j)$, $\lambda_k^G$ and $r^G_{\rho, i,
\varepsilon_1, \dots, \varepsilon_t, \nu_1, \dots, \nu_t}$ (cf.\ Def\/inition~\ref{equivlift} and
Corollary~\ref{eqmodgenerators}), which generate $KR_G^*(G)$ as a~$KR_G^*(\pt)$-algebra, as a~result of
Corollary~\ref{equivstrthmg}.
We show that $(\lambda_k^G)^2=0$ (cf.\ Proposition~\ref{cplxsquare}).
\item
We compute the module structure of $KR^*_{(U(n),
\sigma_\mathbb{F})}(U(n),
\sigma_\mathbb{F})$ for $\mathbb{F}=\mathbb{R}$ and $\mathbb{H}$.
\item
Let $T$ be the maximal torus of diagonal matrices in $U(n)$ and, by abuse of notation, $\sigma_\mathbb{R}$ be the inversion map on $T$, $\sigma_\mathbb{H}$ be the involution on $U(n)/T$ (where $n=2m$ is even) def\/ined by $gT\mapsto J_m\overline{g}T$.
We show that the restriction map
\begin{gather*}
p^*_G: \  KR^*_{(U(n),
\sigma_\mathbb{R})}(U(n),
\sigma_\mathbb{R})\to KR^*_{(T,
\sigma_\mathbb{R})}(T,
\sigma_\mathbb{R})
\end{gather*}
and the map
\begin{gather*}
q^*_G: \ KR^*_{(U(2m),
\sigma_\mathbb{H})}(U(2m),
\sigma_\mathbb{H})\to KR^*_{(U(2m),
\sigma_\mathbb{H})}(U(2m)/T\times T,
\sigma_\mathbb{H}
\times
\sigma_\mathbb{R})
\end{gather*}
induced by the Weyl covering map $q_G: U(2m)/T\times T\to U(2m)$, $(gT, t)\mapsto gtg^{-1}$, are injective.
\item
Let $\sigma_n$ be the class of the standard representation of $U(n)$.
We pass the computation of the two squares $\delta_\mathbb{R}^G(\sigma_n)^2\in KR^*_{(U(n), \sigma_\mathbb{R})}(U(n),
\sigma_\mathbb{R})$ and $\delta_\mathbb{H}^G(\sigma_{2m})^2\in KR^*_{(U(2m),
\sigma_\mathbb{H})}(U(2m),
\sigma_\mathbb{H})$ through the induced map $p^*_G$ and $q_G^*$ to their images and get equations~\eqref{realsigmasquare}
and~\eqref{quatsigmasquare} in Proposition~\ref{equivkrthysquaresigma}.
\item
Applying induced maps $\varphi_i^*$ and $\theta_j^*$ to equations~\eqref{realsigmasquare} and~\eqref{quatsigmasquare}
yields equations~\eqref{realsquare} and~\eqref{quatsquare} in Theorem~\ref{equivkrthysquare} which, together with
Proposition~\ref{cplxsquare} and some relations among~$\eta$,~$\mu$ and $r^G_{\rho, i, \varepsilon_1, \dots,
\varepsilon_t, \nu_1, \dots, \nu_t}$ deduced from Proposition~\ref{krprelim}, describe completely the ring structure of
$KR_G^*(G)$ (cf.\ Theorem~\ref{mainthm2}).
\end{enumerate}
\begin{Remark}\quad
\begin{enumerate}\itemsep=0pt
\item
Seymour f\/irst suggested the analogues of Steps 3, 4 and 5 in the ordinary $KR$-theory case in~\cite{Se} in an attempt to
compute $\delta_\mathbb{R}(\varphi_i)^2$ and $\delta_\mathbb{H}(\theta_j)^2$, but failed to establish Step~3, which he
assumed to be true to make conjectures about $\delta_\mathbb{R}(\varphi_i)^2$.
\item
In equivariant complex $K$-theory, $K_G^*(G/T\times T)\cong K_T^*(T)$ for any compact Lie group~$G$, and the two maps
$p_G^*$ (the restriction map induced by the inclusion $T\hookrightarrow G$) and $q_G^*$ which is induced by the Weyl
covering map are the same.
If $\pi_1(G)$ is torsion-free, then these two maps are shown to be injective (cf.~\cite{BZ}.
In fact it is even shown there that the maps inject onto the Weyl invariants of $K_T^*(T)$).
In the case of equivariant $KR$-theory, things are more complicated.
First of all, while in the case where $(G,\sigma_G)=(U(n), \sigma_\mathbb{R})$, it is true that $KR_G^*(G/T\times T)\cong KR_T^*(T)$, and $p_G^*$ and $q_G^*$ are the same, it is no longer true in
the case where $(G,\sigma_G)=(U(2m), \sigma_\mathbb{H})$.
In Step 3, we use $q_G^*$ for the quaternionic type involution case because we f\/ind that it admits an easier description
than $p_G^*$ does.
Second, we do not know whether $p_G^*$ and $q_G^*$ are injective for general compact Real Lie groups (equipped with any
Lie group involution).
For our purpose it is suf\/f\/icient to show the injectivity results in Step 3.
\end{enumerate}
\end{Remark}

\subsection{A structure theorem}
\begin{Definition}
\label{realequivformal}
A~$G$-space $X$ is a~\emph{Real equivariantly formal} space if
\begin{enumerate}\itemsep=0pt
\item[1)]
$G$ is a~compact Real Lie group,
\item[2)]
$X$ is a~weakly equivariantly formal~$G$-space, and
\item[3)]
the forgetful map $KR^*_G(X)\to KR^*(X)$ admits a~section $s_R: KR^*(X)\to KR^*_G(X)$ which is
a~$KR^*(\pt)$-module homomorphism.
\end{enumerate}
\end{Definition}
\begin{Remark}
If $X$ is a~weakly equivariantly formal~$G$-space, then the forgetful map $K_G^*(X)\to K^*(X)$ admits a~(not necessarily
unique) section $s: K^*(X)\to K_G^*(X)$ which is a~group homomorphism.
\end{Remark}
\begin{Definition}
For a~section $s: K^*(X)\to K^*_G(X)$ (resp.\
$s_R: KR^*(X)\to KR_G^*(X))$ and $a\in K^*(X)$ (resp.\
$a\in KR^*(X)$), we call $s(a)$ (resp.~$s_R(a)$) a~\emph{$($Real$)$ equivariant lift} of~$a$, with respect to $s$ (resp.~$s_R$).
\end{Definition}

We f\/irst prove a~structure theorem of equivariant $KR$-theory of Real equivariantly formal spaces.
\begin{Theorem}
\label{equivstrthm}
Let $X$ be a~Real equivariantly formal space.
For any element $a\in K^*(X)$ $($resp.\
$a\in KR^*(X))$, let $a_G\in K_G^*(X)$ $($resp.\
$a_G\in KR_G^*(X))$ be a~$($Real$)$ equivariant lift of $a$ with respect to a~group homomorphic section $s$ $($resp.~$s_R$ which is a~$KR^*(\pt)$-module homomorphism$)$.
Then the map
\begin{gather*}
\begin{split}
& f:  \ (RR(G, \mathbb{R})\oplus RH(G, \mathbb{R}))\otimes KR^*(X)\oplus r(R(G, \mathbb{C})\otimes K^*(X))  \to  KR_G^*(X),
\\
& \phantom{f: {}} \
\rho_1\otimes a_1\oplus r(\rho_2\otimes a_2)  \mapsto  \rho_1\cdot (a_1)_{G}\oplus r(\rho_2\cdot (a_2)_G).
\end{split}
\end{gather*}
is a~group isomorphism.
In particular, if $R(G, \mathbb{C})=0$, then $f$ is a~$KR_G^*(\pt)$-module isomorphism.
\end{Theorem}

\begin{proof}
Consider the following $H(p, q)$-systems
\begin{gather*}
HR^\alpha(p, q):=KR^{-\alpha}\big(X\times S^{q, 0}, X\times S^{p, 0}\big)\cong KR^{-\alpha+p}\big(X\times S^{q-p, 0}\big),
\\
H^{\alpha}(p, q):= K^{-\alpha}\big(X\times S^{q, 0}, X\times S^{p, 0}\big)\cong K^{-\alpha+p}\big(X\times S^{q-p, 0}\big),
\\
HR_G^{\alpha}(p, q):=KR_G^{-\alpha}\big(X\times S^{q, 0}, X\times S^{p, 0}\big)\cong KR_G^{-\alpha+p}\big(X\times S^{q-p, 0}\big).
\end{gather*}
For the last $H(p, q)$-system,~$G$ acts on $S^{q-p, 0}$ trivially.
The spectral sequences induced by these $H(p, q)$-systems converge to $KR^*(X)$, $K^*(X)$ and $KR_G^*(X)$ respectively
(for the assertion for the f\/irst two $H(p, q)$-systems,
see the proofs of Theorem~3.1 and Lemma~4.1 of~\cite{Se}.
That the third $H(p, q)$-system converges to $KR_G^*(X)$ follows from a~straightforward generalization of the
aforementioned proofs by adding equivariant structure throughout).
Consider the two long exact sequences for the pair $(X\times B^{q-p, 0}, X\times S^{q-p, 0})$, with the top exact
sequence involving equivariant $KR$-theory and the bottom one ordinary $KR$-theory, and the vertical maps being
forgetful maps.
By applying the f\/ive-lemma, we have that each element in the f\/irst two $H(p, q)$-systems has a~(Real) equivariant lift.
Def\/ine a~group homomorphism
\begin{gather*}
f(p, q): \ (RR(G, \mathbb{R})\oplus RH(G, \mathbb{R}))\otimes HR^\alpha(p, q)\oplus r(R(G, \mathbb{C})\otimes H^\alpha(p,
q))\to HR_G^{\alpha}(p, q)
\end{gather*}
by
\begin{gather*}
\rho_1\otimes a_1\oplus r(\rho_2\otimes a_2)\mapsto \rho_1\cdot (a_1)_{G}\oplus r(\rho_2\cdot(a_2)_G).
\end{gather*}
As $RR(G, \mathbb{R})$, $RH(G, \mathbb{R})$ and $R(G, \mathbb{C})$ are free abelian groups, and tensoring free abelian
groups and taking cohomology commute, $f$ is the abutment of $f(p, q)$.
On the $E_1^{p, q}$-page, $f(p, q)$ becomes
\begin{gather*}
f_1^{p, q}: \ (RR(G, \mathbb{R})\oplus RH(G, \mathbb{R}))\otimes KR^{-q}\big(X\times S^{1, 0}\big)\oplus r(R(G, \mathbb{C})\otimes
K^{-q}\big(X\times S^{1, 0}\big))
\\
\hphantom{f_1^{p, q}: {}}{} \ \to
KR_G^{-q}\big(X\times S^{1, 0}\big).
\end{gather*}
Note that $K^{-q}(X\times S^{1, 0})\cong K^{-q}(X)\oplus K^{-q}(X)$, $KR^{-q}(X\times S^{1, 0})\cong K^{-q}(X)$, and
$KR_G^{-q}(X\times S^{1, 0})\cong K_G^{-q}(X)$.
With the above identif\/ication,
\begin{gather*}
r: \ R(G, \mathbb{C})\otimes K^{-q}\big(X\times S^{1, 0}\big)  \to  KR_G^{-q}\big(X\times S^{1, 0}\big)\cong K_G^{-q}(X),
\\
\phantom{r: {}} \
\rho_1\otimes(a_1, 0)\oplus \rho_2\otimes (0, a_2)  \mapsto  \rho_1\cdot (a_1)_G+\overline{\sigma_G^*}
\rho_2\cdot \big(\overline{\sigma_G^*}a_2\big)_G.
\end{gather*}
So $r(R(G, \mathbb{C})\otimes K^{-q}(X\times S^{1, 0}))= R(G, \mathbb{C})\otimes K^{-q}(X)$.
$f_1^{p, q}$ is a~group homomorphism from $(RR(G, \mathbb{R})\oplus RH(G, \mathbb{R})\oplus R(G, \mathbb{C}))\otimes
K^{-q}(X)\cong R(G)\otimes K^{-q}(X)$ to $K_G^{-q}(X)$, which is an isomorphism by weak equivariant formality of $X$.
It follows that $f$ is also an isomorphism.
If $R(G, \mathbb{C})=0$, then by~(1) of Proposition~\ref{KR_G} and~(3) of Def\/inition~\ref{realequivformal}, $f$ is
indeed a~$KR_G^*(\pt)$-module isomorphism.
\end{proof}
\begin{Remark}
The term `Real equivariant formality' is suggested by the observation that, if $X$ is a~Real equivariantly formal
$G$-space and $R(G, \mathbb{C})=0$, then the map
\begin{gather*}
KR^*_G(X)\otimes_{RR(G, \mathbb{R})\oplus RH(G, \mathbb{R})}\mathbb{Z}\to KR^*(X)
\end{gather*}
induced by the forgetful map is a~ring isomorphism, which smacks of the ring isomorphism in the def\/inition of weak
equivariant formality.
\end{Remark}
\begin{Lemma}
\label{realequivlift}
$\delta_\mathbb{R}(\varphi_i)$, $\delta_\mathbb{H}(\theta_j)$, $\lambda_k$ and $r_{i, \varepsilon_1, \dots,
\varepsilon_t, \nu_1, \dots, \nu_t}\in KR^*(G)$ all have Real equivariant lifts in $KR_G^*(G)$.
Hence~$G$ is a~Real equivariantly formal space.
\end{Lemma}
\begin{proof}
A natural choice of a~Real equivariant lift of $\delta_\mathbb{R}(\varphi_i)$ is represented by the complex of vector
bundles in Proposition~\ref{-1_class} equipped with both the Real structure and equivariant structure def\/ined for
$\delta_\mathbb{R}(\varphi_i)$ and $\delta_G(\varphi)$ respectively.
These two structures are easily seen to be compatible.
A Real equivariant lift of $\delta_\mathbb{H}(\theta_j)$ can be similarly def\/ined.
The class
\begin{gather*}
r\big(\beta^i\cdot\delta_G(\gamma_1)^{\varepsilon_1}\cdots\delta_G(\gamma_t)^{\varepsilon_t}\delta_G(\overline{\sigma_G^*}
\gamma_1)^{\nu_1}\cdots\delta_G(\overline{\sigma_G^*}
\gamma_t)^{\nu_t}\big)
\end{gather*}
obviously is a~Real equivariant lift of $r_{i, \varepsilon_1, \dots, \varepsilon_t, \nu_1, \dots, \nu_t}$.
By adding the natural equivariant structure throughout the construction of $\lambda_k$ in the proof of Proposition~4.6
in~\cite{Se}, one can obtain a~Real equivariant lift of $\lambda_k$.
\end{proof}

\begin{Definition}
\label{equivlift}
We f\/ix a~choice of equivariant lift of any element $a\in K^*(G)$ by def\/ining $\delta_G(\rho)$ to be the equivariant lift
of $\delta(\rho)$.
Similarly, we f\/ix a~choice of Real equivariant lift of $a\in KR^*(G)$ by def\/ining $\delta_\mathbb{R}^G(\varphi_i)$,
$\delta_\mathbb{H}^G(\theta_j)$, $\lambda_k^G$, and $r^G_{i, \varepsilon_1, \dots, \varepsilon_t, \nu_1, \dots, \nu_t}$
in the proof of Lemma~\ref{realequivlift} to be the Real equivariant lift of $\delta_\mathbb{R}(\varphi_i)$,
$\delta_\mathbb{H}(\theta_j)$, $\lambda_k$ and $r_{i, \varepsilon_1, \dots, \varepsilon_t, \nu_1, \dots, \nu_t}$.
\end{Definition}
\begin{Remark}
\label{cplxeqsquare}
$\lambda^G_k$ satisf\/ies $c(\lambda^G_k)=\beta^3\delta_G(\gamma_k)\delta_G(\overline{\sigma_G^*}
\gamma_k)$.
\end{Remark}
\begin{Corollary}
\label{equivstrthmg}
Let~$G$ be a~compact, connected and simply-connected Real Lie group.
The map
\begin{gather*}
f: \ (RR(G, \mathbb{R})\oplus RH(G, \mathbb{R}))\otimes KR^*(G)\oplus r(R(G, \mathbb{C})\otimes K^*(G))  \to  KR_G^*(G),
\\
\phantom{f: {}} \
\rho_1\otimes a_1\oplus r(\rho_2\otimes a_2)  \mapsto  \rho_1\cdot (a_1)_{G}\oplus r(\rho_2\cdot (a_2)_G)
\end{gather*}
is a~group isomorphism.
Here $(a_i)_G$ is the $($Real$)$ equivariant lift defined as in Definition~{\rm \ref{equivlift}}.
In particular, if $R(G, \mathbb{C})=0$, then $f$ is an isomorphism of $KR^*_{G}(\pt)$-modules from
$KR_G^*(\pt)\otimes K^*(G)$ to $KR_G^*(G)$.
\end{Corollary}
\begin{proof}
The result follows from Theorem~\ref{equivstrthm} and Lemma~\ref{realequivlift}.
In the special case where $R(G, \mathbb{C})=0$, $KR^*(G)$ is isomorphic to $KR^*(\pt)\otimes K^*(G)$ as
$KR^*(\pt)$-modules by (1) of Theorem~\ref{krtheorymodofg}, and applying Theorem~\ref{equivstrthm} and
Proposition~\ref{KR_G} give $KR^*_G(G)\cong RR(G)\otimes KR^*(G)\cong RR(G)\otimes KR^*(\pt)\otimes K^*(G)\cong
KR^*_G(\pt)\otimes K^*(G)$.
In this way $f$ is a~$KR_G^*(\pt)$-module isomorphism from $KR_G^*(\pt)\otimes K^*(G)$ to $KR_G^*(G)$.
\end{proof}
\begin{Corollary}
\label{eqmodgenerators}
Let
\begin{gather*}
r^G_{\rho, i, \varepsilon_1, \dots, \varepsilon_t, \nu_1, \dots,\nu_t}
:=r\big(\beta^i\cdot\rho\delta_G(\gamma_1)^{\varepsilon_1}\cdots\delta_G(\gamma_t)^{\varepsilon_t}\delta_G(\overline{\sigma_G^*}
\gamma_1)^{\nu_1}\cdots\delta_G(\overline{\sigma_G^*}
\gamma_t)^{\nu_t}\big),
\end{gather*}
where $\rho\in R(G, \mathbb{C})\oplus \mathbb{Z}\cdot\rho_{\text{\rm triv}}$ and $\varepsilon_1, \dots, \varepsilon_t$,
$\nu_1, \dots, \nu_t$ are as in Theorem~{\rm \ref{strthm}}.
Then $KR_G^*(G)$, as an algebra over $KR_G^*(\pt)$, is generated by $\delta_\mathbb{R}^G(\varphi_1), \dots,
\delta_\mathbb{R}^G(\varphi_r)$, $\delta_\mathbb{H}^G(\theta_1), \dots, \delta_\mathbb{H}^G(\theta_s)$, $\lambda_1, \dots,
\lambda_t$, and $r^G_{\rho, i, \varepsilon_1, \dots, \varepsilon_t, \nu_1, \dots, \nu_t}$.
\end{Corollary}
\begin{Remark}
If $\rho=\rho_{\text{triv}}$, then $r^G_{\rho, i, \varepsilon_1, \dots, \varepsilon_t, \nu_1, \dots, \nu_t}=r^G_{i,
\varepsilon_1, \dots, \varepsilon_t, \nu_1, \dots, \nu_t}$.
If $\rho\in R(G, \mathbb{C})$, then $r^G_{\rho, i, \varepsilon_1, \dots, \varepsilon_t, \nu_1, \dots, \nu_t}$ comes from
$r(R(G, \mathbb{C})\otimes K^*(G))$ in the decomposition of Theorem~\ref{equivstrthm}.
\end{Remark}

Now we are in a~position to compute $(\lambda_k^G)^2$ by imitating the proof of Proposition~4.7 in~\cite{Se}.
\begin{Proposition}
\label{cplxsquare}
$\big(\lambda_k^G\big)^2=0$.
\end{Proposition}
\begin{proof}
Consider the Real Lie group $(U(n)\times U(n),
\sigma_\mathbb{C})$, where $\sigma_\mathbb{C}
(g_1, g_2)=(\overline{g_2}, \overline{g_1})$.
Let $p_j: U(n)\times U(n)\to U(n)$ be the projection onto the $j$-th factor, and $u_i=p_1^*(\wedge^i
\sigma_n)$, $v_i=p_2^*(\wedge^i\sigma_n)$.
Thus $\overline{\sigma_\mathbb{C}^*}u_i=v_i$.
A decomposition of $K^*(U(n)\times U(n))$ by the induced involution $\overline{\sigma_\mathbb{C}^*}$ is given by $M\oplus T\oplus \overline{\sigma_\mathbb{C}^*}T$, where $M$ is the subalgebra generated by $\delta(u_1)\delta(v_1), \dots, \delta(u_n)\delta(v_n)$.
By Proposition~\ref{strthm}, there exist $h_1, \dots, h_n\in KR^0(U(n)\times U(n),
\sigma_\mathbb{C})$ such that $c(h_i)=\beta^3\delta(u_i)\delta(v_i)$, and $KR^*(U(n)\times U(n),
\sigma_\mathbb{C})\cong F\oplus r(K^*(+)\otimes T)$, where $F$ is the $KR^*(\pt)$-module freely generated by monomials in $h_1,
\dots, h_n$.
By Corollary~\ref{equivstrthmg},
\begin{gather*}
KR^*_{(U(n)\times U(n),
\sigma_\mathbb{C})}(U(n)\times U(n),
\sigma_\mathbb{C})
\cong
RR(U(n)\times U(n),
\sigma_\mathbb{C}, \mathbb{R})\otimes(F\oplus r(K^*(+)\otimes T))
\\
\qquad{}
\oplus
r(R(U(n)\times U(n),
\sigma_\mathbb{C}, \mathbb{C})\otimes K^*(U(n)\times U(n))).
\end{gather*}
Let $h_i^G$ be the equivariant lift of $h_i$ as def\/ined in Def\/inition~\ref{equivlift}.
So $c(h_i^G)=\beta^3\delta_G(u_i)\delta_G(v_i)$ and $c((h_i^G)^2)=0$.
Consequently $(h_i^G)^2=\eta k_i$ for some $k_i\in KR^{-7}_{(U(n)\times U(n),
\sigma_\mathbb{C})}(U(n)\times U(n),
\sigma_\mathbb{C})$ (cf.\ Gysin sequence~(3.4) in~\cite{At3} and its equivariant analogue).
Since $\eta\cdot r(\cdot)=0$, we may assume that~$k_i$ is from the component $RR(U(n)\times U(n),
\sigma_\mathbb{C}, \mathbb{R})\otimes F$.
But the degree $-7$ piece of the later is~0.
So $(h_i^G)^2=0$.

Consider the map
\begin{gather*}
\gamma_k\times \overline{\sigma_\mathbb{C}^*}\gamma_k: \ (G,\sigma_G)  \to (U(n)\times U(n), \sigma_\mathbb{C}),
\\
\phantom{\gamma_k\times \overline{\sigma_\mathbb{C}^*}\gamma_k: {}} \
g  \mapsto  (\gamma_k(g), \overline{\gamma_k(\sigma_G(g))}).
\end{gather*}
It can be easily seen that $(\gamma_k\times\overline{\sigma_\mathbb{C}^*}\gamma_k)^*(h_1^G)=\lambda_k^G$.
So $(\lambda_k^G)^2=0$.
\end{proof}

\subsection[The module structure of $KR^*_{(U(n), \sigma_\mathbb{F})}(U(n),\sigma_\mathbb{F})$]{The
module structure of $\boldsymbol{KR^*_{(U(n), \sigma_\mathbb{F})}(U(n),\sigma_\mathbb{F})}$}

\begin{Definition}
Let $\sigma_n$ be (the class of) the standard representation of $U(n)$.
\end{Definition}

\begin{Proposition}
\label{realrepunitary}
$\sigma_n,\wedge^2\sigma_n, \dots, \wedge^n\sigma_n\in RR(U(n), \sigma_\mathbb{R}, \mathbb{R})$, $\wedge^{2i}
\sigma_{2m}
\in RR(U(2m),
\sigma_\mathbb{H}, \mathbb{R})$ and $\wedge^{2i+1}
\sigma_{2m}
\in RH(U(2m),
\sigma_\mathbb{H}, \mathbb{R})$.
Also, both $R(U(n),
\sigma_\mathbb{R}, \mathbb{C})$ and $R(U(2m),
\sigma_\mathbb{H}, \mathbb{C})$ are~$0$.
\end{Proposition}
\begin{proof}
For the involution $\sigma_\mathbb{R}$ and $\wedge^i
\sigma_n$, def\/ine the bilinear form
\begin{gather*}
B_\mathbb{R}: \  \wedge^i\sigma_n\times\sigma_\mathbb{R}^*\wedge^i\sigma_n  \to  \mathbb{C},
\\
\phantom{B_\mathbb{R}: {}} \
(v_1\wedge\dots\wedge v_i, w_1\wedge\dots\wedge w_i)  \mapsto  \det(\langle v_j, \overline{w_k}\rangle).
\end{gather*}
Obviously the form is $U(n)$-invariant, symmetric and non-degenerate.
By Proposition~\ref{formclassification}, each of $\wedge^i
\sigma_n$, $1\leq i\leq n$ is a~Real representation of real type.
Similarly, def\/ine, for the involution $\sigma_\mathbb{H}$ and $\wedge^i
\sigma_{2m}$, a~bilinear form
\begin{gather*}
B_\mathbb{H}: \ \wedge^i\sigma_{2m}\times\sigma_\mathbb{H}^*{\wedge^i\sigma_{2m}}  \to  \mathbb{C},
\\
\phantom{B_\mathbb{H}: {}} \
(v_1\wedge\dots\wedge v_i, w_1\wedge\dots\wedge w_i)  \mapsto  \det(\langle J_mv_j, \overline{w_k}\rangle).
\end{gather*}
It is $U(n)$-invariant because
\begin{gather*}
B_\mathbb{H}(gv, gw)=\det(\langle J_mgv_j, \overline{J_m\overline{g}J_m^{-1}w_k}\rangle)
=\det(\langle J_mgv_j, J_mgJ_m^{-1}\overline{w_k}\rangle)
\\
\phantom{B_\mathbb{H}(gv, gw)}
=\det(\langle v_j, J_m^{-1}\overline{w_k}\rangle)
=\det(\langle J_mv_j, \overline{w_k}\rangle).
\end{gather*}
Moreover
\begin{gather*}
B_\mathbb{H}(v, w)=\det(\langle J_mv_j, \overline{w_k}\rangle)
=\det(\langle -v_j, J_m\overline{w_k}\rangle)
=\det(-\overline{\langle J_m\overline{w_k}, v_j\rangle})
\\
\phantom{B_\mathbb{H}(v, w)}
=\det(-\langle J_mw_k, \overline{v_j}\rangle)
=(-1)^iB_\mathbb{H}(w, v).
\end{gather*}
So by Propositions~\ref{realrepclassification},~\ref{formclassification} and~\ref{quatrepclassification}, $\wedge^i
\sigma_{2m}$ is a~Real representation of real type when $i$ is even and a~Quaternionic representation
of real type when $i$ is odd.
There are no complex representations of complex type because
$\wedge^i\sigma_n\cong \sigma_\mathbb{F}^*\overline{\wedge^i\sigma_n}$ for $\mathbb{F}=\mathbb{R}$ and $\mathbb{H}$.
\end{proof}

\begin{Lemma}
%\label{krunitary}
$KR^*_{(U(n),
\sigma_\mathbb{\mathbb{F}})}\!(U(n){,}
\sigma_\mathbb{F})$ is isomorphic to $\Omega^*_{KR^*_{(U(n),
\sigma_\mathbb{F})}(\pt)/ KR^*(\pt)}\!$ as $KR^*_{(U(n),
\sigma_\mathbb{F})}(\pt)$-modules,
where $\mathbb{F}=\mathbb{R}$ or $\mathbb{H}$.
\end{Lemma}
\begin{proof}
Theorem~\ref{equivstrthm} implies that, $KR^*_{(U(n),
\sigma_\mathbb{R})}(U(n),
\sigma_\mathbb{R})\cong RR(U(n),
\sigma_\mathbb{R})\otimes KR^*(U(n),
\sigma_\mathbb{R})$ and $KR^*_{(U(n),
\sigma_\mathbb{H})}(U(n),
\sigma_\mathbb{H})\cong (RR(U(n),
\sigma_\mathbb{H}, \mathbb{R})\oplus RH(U(n),
\sigma_\mathbb{H}, \mathbb{R}))\otimes KR^*(U(n),
\sigma_\mathbb{H})$.
Moreover, by Theorem~\ref{spkrtheorymodofg} and Proposition~\ref{realrepunitary},
\begin{gather*}
KR^*(U(n),
\sigma_\mathbb{R})\cong \bigwedge\nolimits_{KR^*(\pt)}(\delta_\mathbb{R}(\sigma_n), \dots, \delta_\mathbb{R}
(\wedge^n
\sigma_n))
\end{gather*}
and
\begin{gather*}
KR^*(U(2m),
\sigma_\mathbb{H})\cong\bigwedge\nolimits_{KR^*(\pt)}\big(\delta_\mathbb{H}(\sigma_{2m}), \delta_\mathbb{R}(\wedge^2
\sigma_{2m}), \dots, \delta_\mathbb{R}(\wedge^{2m}\sigma_{2m})\big).
\end{gather*}
Putting all these together and applying Theorem~\ref{KR_G}, we get the desired conclusion.
\end{proof}
\begin{Remark}
\label{ungradediso}
As ungraded $KR^*(\pt)$-modules, both
\[
KR_{(U(2m),
\sigma_\mathbb{R})}^*(U(2m),
\sigma_\mathbb{R}) \qquad \text{and}\qquad KR^*_{(U(2m),
\sigma_\mathbb{H})}(U(2m),
\sigma_\mathbb{H})
\]
 are isomorphic to $K^*_{U(2m)}(U(2m))\otimes KR^*(\pt)$.
\end{Remark}

\subsection{Injectivity results}
This step involves proving that the restriction map $p_G^*$ to the equivariant $KR$-theory of the maximal torus and the
map $q_G^*$ induced by the Weyl covering map are injective.
\begin{Lemma}
\label{ringinj}
Let~$G$ be a~compact Lie group and $X$ a~$G$-space.
Let $i_1^*: K_G^*(X)\to K^*_T(X)$ be the map which restricts the~$G$-action to $T$-action.
Then
\begin{gather*}
i_1^*\otimes\Id_R: \  K^*_G(X)\otimes R\to K_T^*(X)\otimes R
\end{gather*}
is injective for any ring $R$.
\end{Lemma}
\begin{proof}
By~\cite[Proposition 4.9]{At4}, $i_1^*$ is split injective.
So is $i_1^*\otimes\Id_R$ for any ring $R$.
\end{proof}
\begin{Lemma}
\label{decomposition}
Let $i_2^*: K_T^*(G)\to K_T^*(T)\cong R(T)\otimes K^*(T)$ be the map induced by the inclusion $T\hookrightarrow G$.
Then
\begin{gather*}
i_2^*\delta_T(\rho)=
\sum\limits_{j=1}^{\dim\rho}e^{\tau_j}\otimes\delta(\tau_j)\in K_T^{-1}(T),
\end{gather*}
where $\tau_j$ are the weights of $\rho$.
\end{Lemma}
\begin{proof}
Let~$V$ be the vector space underlying the representation $\rho$.
$\delta_T(\rho)$ is represented by the complex of $T$-equivariant vector bundles
\begin{gather*}
0\longrightarrow G\times\mathbb{R}\times V \longrightarrow G\times\mathbb{R}\times V\longrightarrow 0,
\\
(g, t, v)  \mapsto  (g, t, -t\rho(g)v) \qquad\text{if} \ \ t\geq 0,
\\
(g, t, v)  \mapsto  (g, t, tv) \qquad\text{if} \ \ t\leq 0
\end{gather*}
which, on restricting to
\begin{gather*}
0\longrightarrow T\times\mathbb{R}\times V
\stackrel{}
{\longrightarrow}T\times\mathbb{R}\times V\longrightarrow 0
\end{gather*}
is decomposed into a~direct sum of complexes of 1-dimensional $T$-equivariant vector bundles, each of which corresponds
to a~weight of $\rho$.
\end{proof}

\begin{Lemma}
\label{Atiyahindex}
Let~$G$ be a~simply-connected, connected compact Lie group and $\rho_1, \dots, \rho_l$ be its fundamental
representations.
Then
\begin{gather}
\label{jacobian}
i_2^*\left(\prod\limits_{i=1}^l\delta_T(\rho_i)\right)=d_G\otimes\prod\limits_{i=1}^l\delta(\varpi_i),
\end{gather}
where $\varpi_i$ the $i$-th fundamental weight and $d_G=
\sum\limits_{w\in W}
\operatorname{sgn}(w)e^{w\cdot
\sum\limits_{i=1}^l\varpi_i}\in R(T)$ is the Weyl deno\-mi\-nator.
\end{Lemma}
\begin{proof}
Equation~\eqref{jacobian} follows from Lemma~\ref{decomposition} and Lemma 3 of~\cite{At2}.
\end{proof}
\begin{Lemma}
\label{wedgeinj}
Let $M$ be an $R$-module freely generated by $m_1, \dots, m_l$, and $N$ an $R$-module.
If $f: M\to N$ is an $R$-module homomorphism, and $rf(m_1)\wedge\dots\wedge f(m_l)\in \bigwedge_R^l N$ is nonzero for
all $r\in R
\setminus\{0\}$, then
\begin{gather*}
\bigwedge\nolimits^* f: \ \bigwedge\nolimits_R^*M\to \bigwedge\nolimits_R^*N
\end{gather*}
is injective.
\end{Lemma}
\begin{proof}
It suf\/f\/ices to show that $\bigwedge^k f$ is injective for $1\leq k\leq l$.
Suppose $I
\subseteq\{1, \dots, l\}$, $|I|=k$, $m_I:=\bigwedge_{i\in I}m_i$ and $f(m_I):=\bigwedge_{i\in I}f(m_i)$.
If $\sum\limits_{|I|=k}
r_Im_I\in\text{ker}(\bigwedge^k f)$, then for any $J$ with $|J|=k$,
\begin{gather*}
0=
\sum\limits_{|I|=k}
r_If(m_I)\wedge f(m_{J^c})=r_Jf(m_1)\wedge\dots\wedge f(m_l).
\end{gather*}
Hence $\sum\limits_{|I|=k}
r_Im_I=0$ and the conclusion follows.
\end{proof}
\begin{Lemma}
\label{twotorsioninj}
Let~$G$ be a~simply-connected, connected and compact Lie group.
Then the map
\begin{gather*}
i_2^*\otimes\Id_R: \ K_T^*(G)\otimes R\to K_T^*(T)\otimes R
\end{gather*}
is injective for any ring $R$.
\end{Lemma}
\begin{proof}
Note that
\begin{gather*}
K_G^*(G)\otimes_{R(G)}R(T)  \to  K_T^*(G),
\\
a\otimes\rho  \mapsto  i_1^*(a)\cdot\rho.
\end{gather*}
is an $R(T)$-algebra isomorphism (cf.~\cite[Theorem~4.4]{HLS}).
Using Theorem~\ref{eqderivation}, we have that $K_T^*(G)$ is isomorphic, as an $R(T)$-algebra, to $\bigwedge_{R(T)}^*M$,
where $M$ is the $R(T)$-module freely generated by $\delta_T(\rho_1), \dots, \delta_T(\rho_l)$.
We also observe that $K_T^*(T)$ is isomorphic, as an $R(T)$-algebra, to $\bigwedge_{R(T)}^*N$, where $N$ is the
$R(T)$-module freely generated by $\delta(\varpi_1), \dots, \delta(\varpi_l)$.
Note that the hypotheses of Lemma~\ref{wedgeinj} are satisf\/ied by $f=i_2^*\otimes\Id_{\mathbb{Z}_m}$ for any
$m\geq 2$, as
$r\prod\limits_{i=1}^li_2^*\delta_T(\rho_i)=ri_2^*\prod\limits_{i=1}^l\delta_T(\rho_i)=rd_G\otimes\prod\limits_{i=1}^l\delta(\varpi_i)$
(by Lemma~\ref{Atiyahindex}) is indeed nonzero for any nonzero $r$ in $\mathbb{Z}_m$ (the coef\/f\/icients of $d_G$ are
either 1 or $-1$, so after reduction mod $m$ $rd_G$ is still nonzero).
Now that $i_2^*\otimes\mathbb{Z}_m$ is injective, so is $i_2^*\otimes\Id_R$ for any ring $R$.
\end{proof}

\begin{Proposition}\label{pginjective}{\samepage\quad
\begin{enumerate}\itemsep=0pt
\item[$1.$]
Let~$G$ be a~compact, connected and simply-connected Real Lie group such that $RR(G, \mathbb{C})=RR(G, \mathbb{H})=0$
and there exists a~maximal torus $T$ on which the involution acts by inversion.
Then the restriction map
$KR_{G}^*(G)\to KR_T^*(T)$
is injective.
\item[$2.$]
The map $p_G^*: KR_{(U(n),
\sigma_\mathbb{R})}^*(U(n),
\sigma_\mathbb{R})\to KR_{(T,
\sigma_\mathbb{R})}^*(T,
\sigma_\mathbb{R})$ is injective.
\end{enumerate}}
\end{Proposition}

\begin{proof}
By Corollary~\ref{equivstrthmg}, $KR^*_G(G)\cong K^*_G(G)\otimes KR^*(\pt)$ and $KR_T^*(T)\cong K^*_T(T)\otimes
KR^*(\pt)$ as $KR^*(\pt)$-modules.
Using this identif\/ication, we can as well identify the restriction map with $i^*\otimes\Id_{KR^*(\pt)}$,
where $i^*:=i_2^*\circ i_1^*$.
Part~(1) then follows from Lemmas~\ref{ringinj} and~\ref{twotorsioninj}.
Part~(2) is immediate once we apply Lemma~\ref{ringinj} and note that the proofs of Lemmas~\ref{Atiyahindex}
and~\ref{twotorsioninj} can be adapted to the case $G=U(n)$ by letting $\sigma_n, \dots, \wedge^n\sigma_n$
play the role of the fundamental representations and their highest weights, the
fundamental weights.
\end{proof}

\begin{Lemma}
\label{quatring}
$KR^*_{(U(2m),
\sigma_\mathbb{H})}(U(2m)/T,
\sigma_\mathbb{H})\cong \mathbb{Z}[e_1^\mathbb{H}, \dots, e_{2m}^\mathbb{H}, (e_1^\mathbb{H}e_2^\mathbb{H}\cdots
e_{2m}^\mathbb{H})^{-1}]\otimes KR^*(\pt)$ as rings, where $e_i^\mathbb{H}$ lives in the degree $-4$ piece.
\end{Lemma}
\begin{proof}
It is known that $K^*(U(2m)/T)\cong\mathbb{Z}[\alpha_1, \dots, \alpha_{2m}]/(s_i-\binom{2m}{i}|1\leq i\leq 2m)$, where
$\alpha_i=[U(2m)\times_T\mathbb{C}_{e_i}]$ and $s_i$ is the $i$-th elementary symmetric polynomial (cf.~\cite[Proposition 2.7.13]{At}).
The induced map $\overline{\sigma_\mathbb{H}^*}$ acts as identity on $K^*(U(2m)/T)$.
The involution $\sigma_\mathbb{H}$ on the base lifts to a~Quaternionic structure
on the associated complex line bundle, so there \mbox{exist}
$\alpha_1^\mathbb{H}, \dots, \alpha_{2m}^\mathbb{H}\in KR^{-4}(U(2m)/T)$, such that their complexif\/ications are
$\beta^2\alpha_1, \dots, \beta^2\alpha_{2m}\in K^*(U(2m)/T)$.
By Theorem~\ref{strthm}, $KR^*(U(2m)/T,
\sigma_\mathbb{H})$ is a~$KR^*(\pt)$-module generated by polynomials in $\alpha_1^\mathbb{H}, \dots, \alpha_{2m}^\mathbb{H}\in
KR^{-4}(U(2m)/T,
\sigma_\mathbb{H})$.
In fact it is not hard to see that $KR^*(U(2m)/T$,
$\sigma_\mathbb{H})$ is isomorphic to
\begin{gather*}
\mathbb{Z}[\alpha_1^\mathbb{H}, \dots, \alpha_{2m}^\mathbb{H}]\otimes
KR^*(\pt)\left/\left.\left(s_{2k}-\binom{2m}{2k}, s_{2k-1}-\frac{1}{2}\mu\binom{2m}{2k-1}\right|1\leq k\leq
m\right)\right..
\end{gather*}
Also obvious is that each of $\alpha_i^\mathbb{H}$ has an equivariant lift $e_i^\mathbb{H}\in KR^{-4}_{(U(2m),
\sigma_\mathbb{H})}(U(2m)/T,\sigma_\mathbb{H})$.

Now that all the hypotheses in Theorem~\ref{equivstrthm} are satisf\/ied, we can apply it, together with the fact that
$R(U(2m),
\sigma_\mathbb{H}, \mathbb{C})=0$ (cf.\ Proposition~\ref{realrepunitary}) to see that $KR^*_{(U(2m),
\sigma_\mathbb{H})}(U(2m)/T,
\sigma_\mathbb{H})$ is isomorphic to
\begin{gather*}
(RR(U(2m),
\sigma_\mathbb{H}, \mathbb{R})\oplus RH(U(2m),
\sigma_\mathbb{H}, \mathbb{R}))\otimes KR^*(U(2m)/T,
\sigma_\mathbb{H})
\end{gather*}
as $RR(U(2m),
\sigma_\mathbb{H}, \mathbb{R})\oplus RH(U(2m),
\sigma_\mathbb{H}, \mathbb{R})$-modules (actually as rings).
Noting that
\begin{gather*}
RR(U(2m),\sigma_\mathbb{H}, \mathbb{R})\oplus RH(U(2m),
\sigma_\mathbb{H}, \mathbb{R})\cong\mathbb{Z}\big[s_1, \dots, s_{2m}, s_{2m}^{-1}\big]
\end{gather*}
we establish the Lemma.
\end{proof}

\begin{Proposition}
%\label{qg}
$KR_{(U(2m),
\sigma_\mathbb{H})}^*(U(2m)/T\times T,
\sigma_\mathbb{H}
\times
\sigma_\mathbb{R})$ is isomorphic to
\begin{gather*}
\mathbb{Z}\big[e_1^\mathbb{H}, \dots, e_{2m}^\mathbb{H},
\big(e_1^\mathbb{H}\cdots e_{2m}^\mathbb{H}\big)^{-1}\big]
\otimes KR^*(T,\sigma_\mathbb{R})
\end{gather*}
as graded rings.
\end{Proposition}
\begin{proof}
First, by~\cite[Theorem~1]{P}, Proposition~\ref{KR_G} and Lemma~\ref{quatring}, $KR^*_{(U(2m),
\sigma_\mathbb{H})}(U(2m)/T,
\sigma_\mathbb{H})$ is a~free $KR_{(U(2m),
\sigma_\mathbb{H})}^*(\pt)$-module.
The same is also true of $KR_{(U(2m),
\sigma_\mathbb{H})}^*(T,
\sigma_\mathbb{R})$ since, by Theo\-rem~\ref{equivstrthm}, it is isomorphic to $RR(U(2m),
\sigma_\mathbb{H})\otimes KR^*(T,
\sigma_\mathbb{R})$, which in turn is isomorphic to $KR^*_{(U(2m),
\sigma_\mathbb{H})}(\pt)\otimes K^*(T)$.
The proposition
follows from a~version of K\"unneth formula for equivariant $KR$-theory.
\end{proof}
\begin{Remark}
\label{ungradedisoh}
$KR^*_{(U(2m),
\sigma_\mathbb{H})}(U(2m)/T\times T,
\sigma_\mathbb{H}
\times
\sigma_\mathbb{R})$ is isomorphic to $K^*_T(T)\otimes KR^*(\pt)$ and $KR^*_{(T,
\sigma_\mathbb{R})}(T,
\sigma_\mathbb{R})$ as ungraded $KR^*(\pt)$-modules.
\end{Remark}
\begin{Proposition}
\label{qginjective}
The map $q_G^*$ is injective.
\end{Proposition}
\begin{proof}
By Remarks~\ref{ungradediso} and~\ref{ungradedisoh}, $KR^*_{(U(2m),
\sigma_\mathbb{H})}(U(2m),
\sigma_\mathbb{H})$ and $KR^*_{(U(2m),
\sigma_\mathbb{H})}(U(2m)/T\times T,
\sigma_\mathbb{H}
\times
\sigma_\mathbb{R})$ are isomorphic to $KR^*_{(U(2m),
\sigma_\mathbb{R})}(U(2m),
\sigma_\mathbb{R})$ and $KR^*_{(T,
\sigma_\mathbb{R})}(T,
\sigma_\mathbb{R})$ respectively, as ungraded $KR^*(\pt)$-modules.
It is not hard to see that $q_G^*$ can be identif\/ied with $p_G^*$ under these isomorphisms.
Now the result follows from Proposition~\ref{pginjective}.
\end{proof}
\subsection{Squares of algebra generators of real and quaternionic types}
\begin{Lemma}
\label{eqkrthytorus}
$KR^*(T,
\sigma_\mathbb{R})$ is isomorphic to the exterior algebra over $KR^*(\pt)$ generated by $\delta_\mathbb{R}(e_1), \dots,
\delta_\mathbb{R}(e_n)$, as $KR^*_{(T,
\sigma_\mathbb{R})}(\pt)$-modules.
Here $e_i$ is the $1$-dimensional complex representation with weight being the $i$-th standard basis vector of the weight
lattice.
Moreover, $\delta_\mathbb{R}(e_i)^2=\eta\delta_\mathbb{R}(e_i)$.
\end{Lemma}

\begin{proof}
Since $R(T,
\sigma_\mathbb{R}, \mathbb{C})=0$, the module structure follows from Theorem~\ref{spkrtheorymodofg}.
For the second part of the Lemma, see the appendix of~\cite{Se}.
\end{proof}
\begin{Proposition}
\label{equivkrthysquaresigma}
In $KR^*_{(U(n),
\sigma_\mathbb{R})}(U(n),
\sigma_\mathbb{R})$
\begin{gather}
\label{realsigmasquare}
\delta_\mathbb{R}^G(\sigma_n)^2
=\eta\big(\sigma_n\cdot\delta_\mathbb{R}^G(\sigma_n)-\delta_\mathbb{R}^G(\wedge^2\sigma_n)\big).
\end{gather}
In $KR^*_{(U(2m),
\sigma_\mathbb{H})}(U(2m),
\sigma_\mathbb{H})$,
\begin{gather}
\label{quatsigmasquare}
\delta_\mathbb{H}^G(\sigma_{2m})^2=\eta\big(\sigma_{2m}
\cdot\delta_\mathbb{H}^G(\sigma_{2m})-\delta_\mathbb{R}^G(\wedge^2\sigma_{2m})\big).
\end{gather}
\end{Proposition}
\begin{proof}
Now that we have shown that $p^*_G$ and $q^*_G$ are injective by Propositions~\ref{pginjective} and~\ref{qginjective},
we can compute $\delta_\mathbb{R}^G(\sigma_n)^2$ and $\delta_\mathbb{H}^G(\sigma_{2m})^2$
by passing the computation through $p^*_G$ and $q_G^*$ to their images.
We prove the case $\mathbb{F}=\mathbb{R}$.
The proof of the case $\mathbb{F}=\mathbb{H}$ is similar so we leave it to the reader.
Note that
\begin{gather*}
p^*_G\big(\delta_\mathbb{R}^G(\sigma_n)^2\big)=p^*_G(\delta_\mathbb{R}^G(\sigma_n))^2
=\left(\sum\limits_{i=1}^n e_i\otimes\delta_\mathbb{R}(e_i)\right)^2
\\
\phantom{p^*_G\big(\delta_\mathbb{R}^G(\sigma_n)^2\big)}=
\sum\limits_{i=1}^n e_i^2\otimes\eta\delta_\mathbb{R}(e_i)+
\sum\limits_{i\neq j}
e_ie_j\otimes\delta_\mathbb{R}(e_i)\delta_\mathbb{R}(e_j)
%\qquad (\text{Lemma }\ref{eqkrthytorus})
\\
\phantom{p^*_G\big(\delta_\mathbb{R}^G(\sigma_n)^2\big)} \overset{\text{Lemma~\ref{eqkrthytorus}}}{=}
\sum\limits_{i=1}^ne_i^2\otimes\eta \delta_\mathbb{R}(e_i)+
\sum\limits_{i<j}
e_i e_j\otimes\big(\delta_\mathbb{R}(e_i)\delta_\mathbb{R}(e_j)+\delta_\mathbb{R}(e_j)\delta_\mathbb{R}(e_i)\big)
\\
\phantom{p^*_G\big(\delta_\mathbb{R}^G(\sigma_n)^2\big)}=
\sum\limits_{i=1}^n e_i^2\otimes\eta\delta_\mathbb{R}(e_i).
\end{gather*}
On the other hand,
\begin{gather*}
p^*(\sigma_n\cdot\delta_\mathbb{R}^G(\sigma_n))=\left(\sum\limits_{i=1}^ne_i\otimes 1\right)
\left(\sum\limits_{i=1}^n e_i\otimes\delta_\mathbb{R}(e_i)\right)
=
\sum\limits_{1\leq i, j\leq n}
e_ie_j\otimes\delta_\mathbb{R}(e_i)
\\
\phantom{p^*(\sigma_n\cdot\delta_\mathbb{R}^G(\sigma_n))}=
\sum\limits_{i=1}^n e_i^2\otimes \delta_\mathbb{R}(e_i)+
\sum\limits_{i\neq j}
e_j e_i\otimes\delta_\mathbb{R}(e_i)
\\
\phantom{p^*(\sigma_n\cdot\delta_\mathbb{R}^G(\sigma_n))}=
\sum\limits_{i=1}^n e_i^2\otimes\delta_\mathbb{R}(e_i)+
\sum\limits_{1\leq i<j\leq n}
e_j e_i\otimes(\delta_\mathbb{R}(e_i)+\delta_\mathbb{R}(e_j))
\\
\phantom{p^*(\sigma_n\cdot\delta_\mathbb{R}^G(\sigma_n))}=
\sum\limits_{i=1}^n e_i^2\otimes\delta_\mathbb{R}(e_i)+p^*_G\big(\delta_\mathbb{R}^G(\wedge^2 \sigma_n)\big).
\end{gather*}
From the above equations we obtain
\begin{gather*}
p^*_G\big(\delta_\mathbb{R}^G(\sigma_n)^2\big)
=\eta p^*_G(\sigma_n\cdot\delta_\mathbb{R}^G(\sigma_n)-\delta_\mathbb{R}^G(\wedge^2\sigma_n)).
\end{gather*}
By Proposition~\ref{pginjective},
\begin{gather*}
\delta_\mathbb{R}^G(\sigma_n)^2
=\eta\big(\sigma_n\cdot\delta_\mathbb{R}^G(\sigma_n)-\delta_\mathbb{R}^G(\wedge^2\sigma_n)\big). \tag*{\qed}
\end{gather*}
\renewcommand{\qed}{}
\end{proof}
\begin{Theorem}
\label{equivkrthysquare}
Let~$G$ be a~Real compact Lie group.
Then
\begin{gather}
\label{realsquare}
\delta_\mathbb{R}^G(\varphi_i)^2
=\eta\big(\varphi_i\cdot\delta_\mathbb{R}^G(\varphi_i)-\delta_\mathbb{R}^G(\wedge^2\varphi_i)\big),
\\
\label{quatsquare}
\delta_\mathbb{H}^G(\theta_j)^2
=\eta\big(\theta_j\cdot\delta_\mathbb{H}^G(\theta_j)-\delta_\mathbb{R}^G(\wedge^2\theta_j)\big).
\end{gather}
\end{Theorem}
\begin{proof}
The induced map $\varphi_i^*: KR^*_{(U(n),
\sigma_\mathbb{R})}(U(n),
\sigma_\mathbb{R})\to KR^*_{(G,
\sigma_G)}
(G,
\sigma_G)$ sends $\sigma_n$ to $\varphi_i$, and $\delta_\mathbb{R}^G(\sigma_n)$ to $\delta_\mathbb{R}^G(\varphi_i)$.
Likewise, the induced map $\theta_j^*: KR^*_{(U(2m),
\sigma_\mathbb{H})}(U(2m),
\sigma_\mathbb{H})$ sends $\sigma_{2m}$ to $\theta_j$, and $\delta_\mathbb{H}^G(\sigma_{2m})$ to $\delta_\mathbb{H}^G(\theta_j)$.
The result now follows from Proposition~\ref{equivkrthysquaresigma}.
\end{proof}

To further express $\delta_\mathbb{R}^G(\wedge^2\varphi_i)$ and $\delta_\mathbb{R}^G(\wedge^2\theta_j)$ in terms of the
module generators associated with the fundamental representations, we may use the following derivation property of
$\delta_\mathbb{R}^G$ and $\delta_\mathbb{H}^G$.
\begin{Proposition}
\label{eqderivative}
$\delta_\mathbb{R}^G\oplus\delta_\mathbb{H}^G$ is a~derivation of the graded ring $RR(G)\oplus RH(G)$ $($with $RR(G)$ of
degree $0$ and $RH(G)$ of degree $-4)$ taking values in the graded module $KR^{-1}_G(G)\oplus KR^{-5}_G(G)$.
\end{Proposition}
\begin{proof}
We refer the reader to the proof of Proposition~3.1 of~\cite{BZ}
with the def\/inition of $\delta_G(\rho)$ given there
(which is incorrect) replaced by the one in Def\/inition~\ref{defdeltag}.
One just need to simply check that the homotopy $\rho_s$ in the proof for $t\geq 0$ intertwines with both $\sigma_\mathbb{R}$ and $\sigma_\mathbb{H}$.
\end{proof}
\begin{Corollary}
\label{krringunitary}
In $KR^*_{(U(n),
\sigma_\mathbb{R})}(U(n),
\sigma_\mathbb{R})$,
\begin{gather*}
\delta_\mathbb{R}^G\big({\wedge}^k
\sigma_n\big)^2=\eta\sum\limits_{i=1}^{2k}\wedge^{2k-i}
\sigma_n\cdot\delta_\mathbb{R}^G\big({\wedge}^i
\sigma_n\big).
\end{gather*}
In $KR^*_{(U(2m),
\sigma_\mathbb{H})}(U(2m),
\sigma_\mathbb{H})$,
\begin{gather*}
\delta_\mathbb{H}^G\big({\wedge}^{2k-1}\sigma_{2m}\big)^2=\eta\sum\limits_{j=1}^{2k-1}\left(\wedge^{4k-2j-1}
\sigma_{2m}\cdot\delta_\mathbb{H}^G\big({\wedge}^{2j-1}\sigma_{2m}\big)+\wedge^{4k-2j-2}\sigma_{2m}
\cdot\delta_\mathbb{R}^G\big({\wedge}^{2j}\sigma_{2m}\big)\right),
\\
\delta_\mathbb{R}^G\big({\wedge}^{2k}\sigma_{2m}\big)^2=\eta\sum\limits_{j=1}^{2k}\left(\wedge^{4k-2j+1}
\sigma_{2m}\cdot\delta_\mathbb{H}^G\big({\wedge}^{2j-1}\sigma_{2m}\big)+\wedge^{4k-2j}\sigma_{2m}
\cdot\delta_\mathbb{R}^G\big({\wedge}^{2j}\sigma_{2m}\big)\right).
\end{gather*}
\end{Corollary}
\begin{proof}
By the def\/inition,
\begin{gather*}
\big({\wedge}^k
\sigma_n\big)^*\big(\delta_\mathbb{R}^G\big({\wedge}^2
\sigma_{\left(
\begin{array}{@{}c@{}}
n
\\
k
\end{array}\right)
}\big)\big)=\delta_\mathbb{R}^G\big({\wedge}^2\big({\wedge}^k
\sigma_n\big)\big).
\end{gather*}
By Exercise 15.32 of~\cite{FH},
\begin{gather*}
\wedge^2(\wedge^k
\sigma_n)=\bigoplus_i\Gamma_{\varpi_{k-2i+1}
+\varpi_{k+2i-1}}.
\end{gather*}
By Giambelli's formula,
\begin{gather*}
\Gamma_{\varpi_{k-2i+1}+\varpi_{k+2i-1}}=\wedge^{k+2i-1}
\sigma_n\cdot\wedge^{k-2i+1}\sigma_n-\wedge^{k+2i}\sigma_n\cdot\wedge^{k-2i}\sigma_n.
\end{gather*}
By Proposition~\ref{eqderivative}
\begin{gather*}
\delta_\mathbb{R}^G(\Gamma_{\varpi_{k-2i+1}+\varpi_{k+2i-1}})=\wedge^{k+2i-1}
\sigma_n\cdot\delta_\mathbb{R}^G\big({\wedge}^{k-2i+1}\sigma_n\big)+\wedge^{k-2i+1}\sigma_n\cdot\delta_\mathbb{R}^G\big({\wedge}^{k+2i-1}\sigma_n\big)
\\
\hphantom{\delta_\mathbb{R}^G(\Gamma_{\varpi_{k-2i+1}+\varpi_{k+2i-1}})=}
{}-\wedge^{k+2i}
\sigma_n\cdot\delta_\mathbb{R}^G\big({\wedge}^{k-2i}\sigma_n\big)-\wedge^{k-2i}
\sigma_n\cdot\delta_\mathbb{R}^G\big({\wedge}^{k+2i}\sigma_n\big).
\end{gather*}
Now the f\/irst equation is immediate.
The second and third equations can be derived similarly.\looseness=1
\end{proof}

Putting together the previous results yields the following full description of the ring structure of $KR_G^*(G)$.
\begin{Theorem}
\label{mainthm2}
Let~$G$ be a~simply-connected, connected and compact Real Lie group.
Viewing~$G$ as a~Real~$G$-space with adjoint action, we have
\begin{enumerate}\itemsep=0pt
\item[$1.$] {\rm (}Corollary~{\rm \ref{equivstrthmg})} The map
\begin{gather*}
f: \ (RR(G, \mathbb{R})\oplus RH(G, \mathbb{R}))\otimes KR^*(G)\oplus r(R(G, \mathbb{C})\otimes K^*(G))\to KR_G^*(G)
\\
\phantom{f: {}} \
\rho_1\otimes a_1\oplus r(\rho_2\otimes a_2)\mapsto \rho_1\cdot (a_1)_{G}\oplus r(\rho_2\cdot (a_2)_G)
\end{gather*}
is a~group isomorphism.
In particular, if $R(G, \mathbb{C})=0$, then $f$ is an isomorphism of $KR^*_{G}(\pt)$-modules.
\item[$2.$] {\rm (}Corollary~{\rm \ref{eqmodgenerators})}
$KR_{G}(G)$ is generated by $\delta_\mathbb{R}^G(\varphi_1), {\dots}, \delta_\mathbb{R}^G(\varphi_r)$,
$\delta_\mathbb{H}^G(\theta_1), {\dots}, \delta_\mathbb{H}^G(\theta_s)$, $\lambda_1^G, {\dots}, \lambda_t^G\!$ and
$r^G_{\rho, i, \varepsilon_1, \dots, \varepsilon_t, \nu_1, \dots, \nu_t}$ as an algebra over $KR_{G}^*(\pt)$.
Moreover,
%\begin{enumerate}\itemsep=0pt
\begin{gather}
(a)\quad{\rm (}Proposition~{\rm \ref{cplxsquare})} \quad \big(\lambda_k^G\big)^2=0\text{ for all }1\leq k\leq t,\nonumber
\\ \label{cplxrel1}
(b)\quad
(r^G_{\rho, i, \varepsilon_1, \dots, \varepsilon_t, \nu_1, \dots, \nu_t})^2
\\
\phantom{(b)\quad}
=
\begin{cases}
\eta^2(\rho\cdot\overline{\sigma_G^*}
\rho)(\lambda^G_1)^{\omega_1}\cdots(\lambda_t^G)^{\omega_t}
&
\text{if }r^G_{\rho, i, \varepsilon_1, \dots, \varepsilon_t, \nu_1, \dots, \nu_t}
\text{ is of degree} \ -1\text{ or }-5,
\\
\pm\mu(\rho\cdot\overline{\sigma_G^*}
\rho)(\lambda^G_1)^{\omega_1}\cdots(\lambda_t^G)^{\omega_t}
&
\text{if }r^G_{\rho, i, \varepsilon_1, \dots, \varepsilon_t, \nu_1, \dots, \nu_t}
\text{ is of degree} \ -2\text{ or }-6,
\\
0
&\text{otherwise}.
\end{cases}
\nonumber
\end{gather}
The sign can be determined using formulae in $(2)$ of Proposition~{\rm \ref{krprelim}}.
\begin{gather}
\label{cplxrel2}
(c)\quad
r^G_{\rho, i, \varepsilon_1, \dots, \varepsilon_t, \nu_1, \dots, \nu_t}\eta=0,\text{ and }
r^G_{\rho, i, \varepsilon_1,\dots, \varepsilon_t, \nu_1, \dots, \nu_t}\mu
=2r^G_{\rho, i+2, \varepsilon_1, \dots, \varepsilon_t, \nu_1, \dots,\nu_t},
\\
(d)\quad
{\rm (}Proposition~{\rm \ref{equivkrthysquare})}
\quad
\delta_\mathbb{R}^G(\varphi_i)^2
=\eta\big(\varphi_i\cdot\delta_\mathbb{R}^G(\varphi_i)-\delta_\mathbb{R}^G(\wedge^2\varphi_i)\big),
\nonumber
\\
\phantom{(d)\quad}
\delta_\mathbb{H}^G(\theta_j)^2
=\eta\big(\theta_j\cdot\delta_\mathbb{H}^G(\theta_j)-\delta_\mathbb{R}^G(\wedge^2\theta_j)\big).
\nonumber
\end{gather}
One can express $\delta_\mathbb{R}^G(\wedge^2\varphi_i)$ and $\delta_\mathbb{H}^G(\wedge^2\theta_j)$ in terms of the
algebra generators using the derivation property of $\delta_\mathbb{R}^G$ and $\delta_\mathbb{H}^G$
$($cf.\ Proposition~{\rm \ref{eqderivative})}.
\end{enumerate}
\end{Theorem}
\begin{proof}
Only~\eqref{cplxrel1} and~\eqref{cplxrel2} need explanation, but they are just equivariant analogues of
Corollary~\ref{krglooseend} and follow from (2) of Proposition~\ref{krprelim} and Remark~\ref{cplxeqsquare}.
\end{proof}
